\chapter{Spectral forms and level-changing tensor products}
\label{cylev:chapter}

\section{Addition of levels}

Let $k$ be a field of characteristic zero.
Let $J$ be a Heisenberg current with singular product
$J(z)J(w)\sim K(z-w)^{-2}$.
In the ordinary tensor product of two such algebras, the current
$J\otimes\mathbf1+\mathbf1\otimes J$ has level $2K$.
The two self-contractions add, and the mixed contractions vanish.
Thus this formula cannot define a vertex-algebra coproduct at a fixed
nonzero level. It instead gives a map from level $K_1+K_2$ to the
tensor product of levels $K_1$ and $K_2$.

There is a further effect when translation changes the contraction.
Over $k[\lambda,c_1,c_2]/(c_1^2,c_2^2)$, consider
\[
 P_c(z)=\frac{\lambda}{z-c}-\frac{\lambda}{z},
\]
where a shifted denominator means its finite nilpotent expansion.
Each separate field has $P_{c_i}(z)=\lambda c_i z^{-2}$, whereas
\begin{equation}\label{cylev:eq:opening-cross}
 P_{c_1+c_2}(z)-P_{c_1}(z)-P_{c_2}(z)
       =\frac{2\lambda c_1c_2}{z^3}.
\end{equation}
The right side is nonzero in this coefficient ring.
A tensor product receiving a coproduct from the summed level must
retain it. We shall construct the extra contraction and prove its
compatibility with fields, iterated coproducts and module actions.

The construction starts with an ordinary Hopf algebra, a translation
derivation and a bilinear form depending on the separation of two
points. The form determines both exchange and contraction. A Hopf
automorphism preserving it specifies a level. The coproduct then
composes levels in their given order, which need not be commutative.
Formal translation levels recover additive parameters and the rational
quantum affine example. Their comparison with Hall charge splitting
requires an additional map between the two coefficient systems.

\section{Rational fields and a spectral form}

Throughout the algebraic construction, $k$ has characteristic zero,
$R$ is a commutative $k$-algebra, and tensor products are over $R$.
States lie in cohomological degree zero. For an $R$-module $V$, put
\begin{equation}\label{cylev:eq:germs}
 \mathcal G_r(V)=V[[z_1,\ldots,z_r]]
 [z_1^{-1},\ldots,z_r^{-1},(z_i-z_j)^{-1}\mid 1\leq i<j\leq r].
\end{equation}
Each element has a finite common denominator. The bound may depend
on the input states. Products and iterates of fields are compared
as expansions of the same element of this module.

\begin{definition}\label{cylev:def:fields}
A \emph{quantum vertex algebra with rational exchange} consists of
a flat module $V$, an operator $T$, and bilinear fields
$Y(a,z)b\in V((z))$.
Its vacuum $\mathbf1$ spans a free direct summand.
They obey
\[
 Y(\mathbf1,z)=\id,\qquad Y(a,z)\mathbf1=e^{zT}a,\qquad
 Y(Ta,z)=\partial_zY(a,z)=[T,Y(a,z)].
\]
They are weakly associative: a product and its consecutive iterate
agree in their common expansion domain after clearing a finite power
of the collision coordinate.

The exchange is an invertible $R[z,z^{-1}]$-linear map
$\mathcal S(z)$ on $(V\otimes V)[z,z^{-1}]$, fixing both vacuum
tensors. It satisfies
\begin{align}
 [T\otimes1,\mathcal S(z)]&=-\partial_z\mathcal S(z),
       \label{cylev:eq:shift}\\
 \mathcal S_{12}(u)\mathcal S_{13}(u+v)\mathcal S_{23}(v)
 &=\mathcal S_{23}(v)\mathcal S_{13}(u+v)\mathcal S_{12}(u),
       \label{cylev:eq:yb}\\
 \mathcal S_{21}(-z)\mathcal S(z)&=1.
       \label{cylev:eq:unitarity}
\end{align}
Here $\mathcal S_{21}$ conjugates $\mathcal S$ by the flip.
Writing $Y(z)(a\otimes b)=Y(a,z)b$, the field identities are
\begin{align}
 &(z-w)^N Y(b,w)Y(a,z)c\nonumber\\
 &\qquad=(z-w)^N Y(z)(1\otimes Y(w))
       \bigl(\mathcal S(z-w)(a\otimes b)\otimes c\bigr),
       \label{cylev:eq:locality}\\
 &\mathcal S(z)(Y(w)\otimes1)
       =(Y(w)\otimes1)\mathcal S_{23}(z)\mathcal S_{13}(z+w).
       \label{cylev:eq:hexagon}
\end{align}
The exponent $N$ in the first identity depends on $a,b$, not on $c$.
In the second identity $\mathcal S_{13}(z+w)$ is expanded at $w=0$.
\end{definition}

This is the rational, algebraic version of the exchange and hexagon
formalism of Etingof and Kazhdan
\cite[\S\S1.3--1.4]{EtingofKazhdanV}.
The formal version below uses completed tensors and these identities
at each finite precision. The definition itself supplies no coproduct.

Let $H$ be a flat Hopf algebra over $R$, with coproduct $\Delta$,
counit $\varepsilon$ and bijective antipode $\mathsf s$.
Write $\Delta a=a_1\otimes a_2$, with the finite sum understood.
Repeated subscripts denote iterated coproducts.
Let $T$ be a derivation and coderivation with $\varepsilon T=0$.
It commutes with $\mathsf s$. Write $a(z)=e^{zT}a$.

\begin{definition}\label{cylev:def:form}
A \emph{unitary spectral coquasitriangular form} is an $R$-bilinear
map $\rho_z:H\otimes H\to R[z,z^{-1}]$ satisfying
\begin{align}
 \rho_z(1,a)&=\varepsilon(a)=\rho_z(a,1),
       \label{cylev:eq:form-unit}\\
 \rho_z(ab,c)&=\rho_z(a,c_1)\rho_z(b,c_2),
       \label{cylev:eq:form-first}\\
 \rho_z(a,bc)&=\rho_z(a_1,c)\rho_z(a_2,b),
       \label{cylev:eq:form-second}\\
 \rho_z(a_1,b_1)a_2(z)b_2
       &=b_1a_1(z)\rho_z(a_2,b_2),
       \label{cylev:eq:form-commutation}\\
 \rho_z(Ta,b)&=\partial_z\rho_z(a,b)=-\rho_z(a,Tb),
       \label{cylev:eq:form-translation}\\
 \rho_z(a_1,b_1)\rho_{-z}(b_2,a_2)
       &=\varepsilon(a)\varepsilon(b).
       \label{cylev:eq:form-unitarity}
\end{align}
The commutation identity takes values in $H[[z]][z^{-1}]$.
The reversed order in \eqref{cylev:eq:form-second} fixes our convention.
\end{definition}

The two multiplication rules alone do not impose exchange or unitarity.
When $T=0$ and the form is constant, the first four rules are the
coquasitriangular identities and the last is cotriangularity.
For two forms, use convolution
$(f*g)(a,b)=f(a_1,b_1)g(a_2,b_2)$, whose unit is
$\varepsilon\otimes\varepsilon$.

\begin{lemma}\label{cylev:lem:form}
The convolution inverse of $\rho_z$ is
\begin{equation}\label{cylev:eq:form-inverse}
 \bar\rho_z(a,b)=\rho_{-z}(b,a)
       =\rho_z(\mathsf sa,b)=\rho_z(a,\mathsf s^{-1}b).
\end{equation}
Simultaneous application of the antipode preserves the form, and
\begin{equation}\label{cylev:eq:form-yb}
\begin{split}
 &\rho_u(a_1,b_1)\rho_{u+v}(a_2,c_1)\rho_v(b_2,c_2)\\
 &\qquad=\rho_v(b_1,c_1)\rho_{u+v}(a_1,c_2)\rho_u(a_2,b_2).
\end{split}
\end{equation}
The latter is an identity in $R[u,v,u^{-1},v^{-1},(u+v)^{-1}]$.
\end{lemma}
\begin{proof}
Apply the first multiplication rule to $\mathsf s(a_1)a_2$ and
$a_1\mathsf s(a_2)$. The antipode identities make
$\rho_z(\mathsf s-,-)$ a two-sided inverse.
The second multiplication rule, applied to
$b_2\mathsf s^{-1}(b_1)$ and $\mathsf s^{-1}(b_2)b_1$, gives the
other inverse. Uniqueness and \eqref{cylev:eq:form-unitarity} give
\eqref{cylev:eq:form-inverse}. Substituting $\mathsf sb$ then gives
$\rho_z(\mathsf sa,\mathsf sb)=\rho_z(a,b)$.

Taylor expansion gives $\rho_v(a(u),c)=\rho_{u+v}(a,c)$.
Pair the two sides of \eqref{cylev:eq:form-commutation}, at $u$,
with $c$ using $\rho_v$. Expanding their products gives exactly the
two sides of \eqref{cylev:eq:form-yb}, expanded at $u=0$.
The expansion map from the displayed localization is injective,
so the rational identity follows.
\end{proof}

Translation puts the commutation identity at two positions:
\begin{equation}\label{cylev:eq:two-positions}
 \rho_{z-w}(a_1,b_1)a_2(z)b_2(w)
       =b_1(w)a_1(z)\rho_{z-w}(a_2,b_2).
\end{equation}
Before introducing contractions, this already gives a regular example.

\begin{proposition}\label{cylev:prop:regular}
On $H$, the fields $Y^{\mathrm{reg}}(a,z)b=a(z)b$ and exchange
\[
 \mathcal S^{\mathrm{reg}}(z)(a\otimes b)
       =\rho_z(a_1,b_1)a_2\otimes b_2\bar\rho_z(a_3,b_3)
\]
form a quantum vertex algebra. Its Hopf coproduct and counit preserve
fields and exchange, using ordinary tensor fields on the target.
\end{proposition}
\begin{proof}
The derivation rule proves the vacuum, translation and associativity
identities. Convolving \eqref{cylev:eq:two-positions} on the right
with $\bar\rho$ proves locality. The finitely many scalar coefficients
on $a,b$ give its bound independently of the terminal state.

For a convolution-invertible form $f$, put
$\Phi_f(a\otimes b)=f(a_1,b_1)a_2\otimes b_2f^{-1}(a_3,b_3)$.
Coassociativity gives $\Phi_f\Phi_g=\Phi_{g*f}$.
Equations \eqref{cylev:eq:form-yb} and
\eqref{cylev:eq:form-inverse} give Yang--Baxter and unitarity.
Differentiating the scalar factors gives \eqref{cylev:eq:shift}.
The first multiplication rule and its reversed inverse rule
$\bar\rho_z(ab,c)=\bar\rho_z(b,c_1)\bar\rho_z(a,c_2)$ give
the hexagon, with the translated argument producing $z+w$.

The Hopf maps preserve multiplication and translation, hence fields.
In the target exchange
$\mathcal S^{\mathrm{reg}}_{13}\mathcal S^{\mathrm{reg}}_{24}$,
the middle factors $\bar\rho_z(a_3,b_3)\rho_z(a_4,b_4)$ cancel.
The result is $(\Delta\otimes\Delta)\mathcal S^{\mathrm{reg}}$.
Applying both counits cancels all scalar factors.
\end{proof}

\section{Levels, coactions and contractions}

\begin{definition}\label{cylev:def:level}
A \emph{spectral level} is a Hopf automorphism $\eta$ of $H$ commuting with
$T$ and satisfying $\rho_z(\eta a,\eta b)=\rho_z(a,b)$.
Products of levels mean composition: $\eta_1\eta_2=\eta_1\circ\eta_2$.
\end{definition}

The levels form a group. It contains $\mathsf s^2$.
For a nilpotent $c\in R$, it also contains $e^{cT}$: this exponential
is finite, and simultaneous translation leaves the form unchanged.
No commutativity of levels is assumed.

Fix a level $\eta$. On the underlying module $H$, define
\begin{equation}\label{cylev:eq:coaction}
 \delta_\eta(b)=b_2\otimes b_3\mathsf s^{-1}(\eta b_1).
\end{equation}
For a right coaction we use parentheses,
$\delta(b)=b_{(0)}\otimes b_{(1)}$, to distinguish its terms from
Hopf coproduct components.

\begin{lemma}\label{cylev:lem:coaction}
The map $\delta_\eta$ is a counital right $H$-coaction, sends $1$ to
$1\otimes1$, and satisfies
\begin{align}
 \delta_\eta T&=(T\otimes1+1\otimes T)\delta_\eta,
       \label{cylev:eq:coaction-translation}\\
 \delta_\eta(xb)&=x_2b_{(0)}\otimes
          x_3b_{(1)}\mathsf s^{-1}(\eta x_1).
       \label{cylev:eq:coaction-product}
\end{align}
\end{lemma}
\begin{proof}
The two coaction composites are both
\[
 b_3\otimes b_4\mathsf s^{-1}(\eta b_2)
       \otimes b_5\mathsf s^{-1}(\eta b_1).
\]
Applying the last counit gives $b$.
The coderivation rule and commutation of $T$ with $\eta,\mathsf s$
give \eqref{cylev:eq:coaction-translation}.
Expanding $\delta_\eta(xb)$ gives the second formula because
$\mathsf s^{-1}$ reverses products.
\end{proof}

For any right comodule $M$, set
\begin{equation}\label{cylev:eq:annihilation}
 A_z^M(x)m=m_{(0)}\rho_{-z}(m_{(1)},x).
\end{equation}
It is a finite Laurent polynomial for fixed $x,m$.
Write $A_z$ for $M=(H,\delta_\eta)$, and let $L(x)$ denote
left multiplication by $x$.

\begin{lemma}\label{cylev:lem:operators}
On every right comodule,
\begin{align}
 A_z(x)A_z(y)&=A_z(yx),\label{cylev:eq:anti-action}\\
 \rho_{w-z}(y_1,x_1)A_z(x_2)A_w(y_2)
 &=A_w(y_1)A_z(x_1)\rho_{w-z}(y_2,x_2).
       \label{cylev:eq:annihilation-exchange}
\end{align}
For $\delta_\eta$ one further has
\begin{align}
 A_z(x)L(b(w))={}&\rho_{w-z}(b_3,x_1)
    \rho_{w-z}(\mathsf s^{-1}\eta b_1,x_3)
    L(b_2(w))A_z(x_2),\label{cylev:eq:mixed-operators}\\
 A_z(x)1={}&\varepsilon(x)1.\label{cylev:eq:annihilation-vacuum}
\end{align}
\end{lemma}
\begin{proof}
Coassociativity writes the first composite as
$m_{(0)}\rho_{-z}(m_{(1)},x)\rho_{-z}(m_{(2)},y)$.
The second multiplication rule combines the scalar factors with
the product $yx$. Pairing the components of $m,x,y$ in
\eqref{cylev:eq:form-yb}, at positions $0,z,w$, gives the second
identity.

For the mixed identity, use \eqref{cylev:eq:coaction-product} on
$b(w)m$. Its last factor is
$b_3(w)m_{(1)}\mathsf s^{-1}\eta b_1(w)$.
The first multiplication rule separates these three factors.
Translation changes the first and last parameters to $w-z$,
leaving $A_z(x_2)$ in the middle. The vacuum identity follows from
$\delta_\eta(1)=1\otimes1$.
\end{proof}

Creation and annihilation now define the fields:
\begin{equation}\label{cylev:eq:field}
\begin{split}
 Y_\eta(a,z)b
 &=L(a_1(z))A_z(a_2)b\\
 &=a_1(z)b_{(0)}\rho_{-z}(b_{(1)},a_2)\\
 &=a_1(z)b_2\rho_{-z}(b_3,a_2)
                  \rho_{-z}(\mathsf s^{-1}\eta b_1,a_3).
\end{split}
\end{equation}
The finite coproducts give an element of $H((z))$.
Denote this state module with its fields by $V_\eta$.

\begin{theorem}\label{cylev:thm:fields}
The fields \eqref{cylev:eq:field}, vacuum $1$, translation $T$ and
exchange
\begin{equation}\label{cylev:eq:exchange}
 \mathcal S_\eta(z)(a\otimes b)
       =a_{(0)}\otimes b_{(0)}\rho_z(a_{(1)},b_{(1)})
\end{equation}
make $V_\eta$ a quantum vertex algebra with rational exchange.
Every finite ordered field product on a fixed state has a common
germ in \eqref{cylev:eq:germs}. Moreover,
\begin{align}
 \delta_\eta(Y_\eta(a,z)b)
 &=Y_\eta(a_{(0)},z)b_{(0)}\otimes b_{(1)}a_{(1)}(z),
       \label{cylev:eq:covariance}\\
 Y_\eta(a_{(0)},z)b_{(0)}\rho_z(a_{(1)},b_{(1)})
 &=e^{zT}Y_\eta(b,-z)a.
       \label{cylev:eq:skew}
\end{align}
\end{theorem}
\begin{proof}
The unit rule for $\rho$ gives the vacuum and creation identities.
The Hopf counit splits off the vacuum summand as an $R$-module.
Using \eqref{cylev:eq:coaction-translation} gives
\[
 [T,A_z(x)]b=-b_{(0)}\rho_{-z}(Tb_{(1)},x)
       =\partial_z(A_z(x)b)=A_z(Tx)b.
\]
Together with the coderivation rule for $\Delta$, this proves both
translation identities for $Y_\eta$.

Move annihilation past creation using
\eqref{cylev:eq:mixed-operators}. The product becomes
\begin{equation}\label{cylev:eq:normal-product}
\begin{split}
 Y_\eta(a,z)Y_\eta(b,w)
 ={}&\rho_{w-z}(b_3,a_2)
       \rho_{w-z}(\mathsf s^{-1}\eta b_1,a_4)\\
 &\quad\cdot L(a_1(z)b_2(w))A_z(a_3)A_w(b_4).
\end{split}
\end{equation}
Expanding the inner field in the iterate instead gives
\begin{equation}\label{cylev:eq:normal-iterate}
\begin{split}
 Y_\eta(Y_\eta(a,u)b,w)
 ={}&\rho_{-u}(b_4,a_3)
       \rho_{-u}(\mathsf s^{-1}\eta b_1,a_4)\\
 &\quad\cdot L(a_1(u+w)b_2(w))A_w(b_3)A_{w+u}(a_2).
\end{split}
\end{equation}
The annihilation exchange identifies these expressions at $u=z-w$.

On a fixed terminal state the remaining annihilations are finite
Laurent polynomials. The other poles lie on $z=w$, so both expressions
come from one element of $\mathcal G_2(H)$.
Repeatedly moving annihilations to the right proves the same assertion
for any finite number of fields.

For covariance, expand the coaction of \eqref{cylev:eq:field}.
The middle factors are interchanged by
\[
 a_3(z)b_{(1)}\rho_{-z}(b_{(2)},a_4)
       =\rho_{-z}(b_{(1)},a_3)b_{(2)}a_4(z),
\]
which is \eqref{cylev:eq:two-positions} at $0,z$.
The outer factors agree directly, proving \eqref{cylev:eq:covariance}.

For skew-symmetry, expand the two coactions on its left, separate
products inside each form, and cancel inverse pairings.
Invariance under $\eta$ leaves
\[
 a_3(z)b_2\rho_z(a_2,b_1)
                  \rho_z(\mathsf s^{-1}\eta a_1,b_3).
\]
Direct expansion of the right side gives
$b_1a_2(z)\rho_z(a_3,b_2)\rho_z(\mathsf s^{-1}\eta a_1,b_3)$.
The commutation identity for $\rho$ equates them.

Associativity and skew-symmetry identify the two orders of a pair
of fields in their collision expansion.
The scalar coefficients in \eqref{cylev:eq:normal-product} have a
pole bound depending only on $a,b$.
The exchanged side has finitely many analogous coefficients.
Their maximum clears the collision denominators independently of
the terminal state, proving \eqref{cylev:eq:locality}.

Coassociativity of $\delta_\eta$ turns
\eqref{cylev:eq:form-yb} and \eqref{cylev:eq:form-unitarity} into
Yang--Baxter and unitarity for $\mathcal S_\eta$.
Differentiation of the paired coaction components gives
\eqref{cylev:eq:shift}.
Finally, the tail of a field in \eqref{cylev:eq:covariance} is
$b_{(1)}a_{(1)}(w)$.
Pairing it with $c_{(1)}$ gives
$\rho_z(b_{(1)},c_{(1)})\rho_{z+w}(a_{(1)},c_{(2)})$,
the two factors of the hexagon.
The unit rule fixes its vacuum tensors.
\end{proof}

\section{The tensor product and its cross contraction}

An \emph{$H$-covariant field algebra} is a quantum vertex algebra
with rational exchange and a right coaction preserving its vacuum
and total translation, satisfying \eqref{cylev:eq:covariance}.
Its exchange is given by \eqref{cylev:eq:exchange}, and every finite
ordered product on a fixed state has a common rational germ.
Morphisms preserve fields, vacuum, translation and coaction.
The algebras $V_\eta$ and the scalar algebra $R$ with trivial coaction
are examples.

For any such algebra $V$, covariance gives
\begin{equation}\label{cylev:eq:annihilation-field}
 A_z^V(x)Y_V(a,w)
       =\rho_{w-z}(a_{(1)},x_2)Y_V(a_{(0)},w)A_z^V(x_1).
\end{equation}
Indeed, its left side pairs $v_{(1)}a_{(1)}(w)$ with $x$.
The first multiplication rule separates precisely the two factors
on the right.

On $V_1\otimes V_2$ define
\begin{align}
 Y_{12}(a\otimes b,z)(c\otimes d)
 &=Y_1(a,z)c_{(0)}\otimes Y_2(b_{(0)},z)d\,
        \rho_{-z}(c_{(1)},b_{(1)}),\label{cylev:eq:tensor-field}\\
 \delta_{12}(a\otimes b)
 &=a_{(0)}\otimes b_{(0)}\otimes b_{(1)}a_{(1)}.
       \label{cylev:eq:tensor-coaction}
\end{align}
The vacuum is $\mathbf1\otimes\mathbf1$ and translation is $T_1+T_2$.
Pairing the last factors of the coaction defines its exchange.
Write $V_1\boxtimes_\rho V_2$ for this field algebra on the ordinary
tensor module. The symbol records its cross contraction at one point.

\begin{theorem}\label{cylev:thm:tensor}
These formulas define a monoidal product on $H$-covariant field
algebras, with unit $R$ and the usual associativity map of tensor
modules.
\end{theorem}
\begin{proof}
The reversed product of tails is a coaction because $\Delta$ is
multiplicative. The coderivation identity gives its translation rule.
The two counits give the vacuum and creation identities.
Differentiating the pairing proves translation for the fields.

In operator notation the tensor field is
\[
 Y_{12}(a\otimes b,z)
       =Y_1(a,z)A_z^{V_1}(b_{(1)})\otimes Y_2(b_{(0)},z).
\]
Using \eqref{cylev:eq:annihilation-field}, a product of two such
fields is
\begin{equation}\label{cylev:eq:tensor-normal}
\begin{split}
 &\rho_{w-z}(c_{(1)},b_{(2)})
   Y_1(a,z)Y_1(c_{(0)},w)
   A_z^{V_1}(b_{(1)})A_w^{V_1}(d_{(1)})\\
 &\hspace{34mm}\otimes Y_2(b_{(0)},z)Y_2(d_{(0)},w).
\end{split}
\end{equation}
For the iterate, covariance in $V_2$ instead gives the annihilation
$A_w^{V_1}(d_{(1)}b_{(1)}(z-w))$.
The anti-action and translation rules turn this into
$A_z^{V_1}(b_{(1)})A_w^{V_1}(d_{(1)})$.
Associativity in the two factors therefore gives
\eqref{cylev:eq:tensor-normal} again.

Finite coactions and the common-germ property in each factor bound
all denominators. Repeating the calculation gives the common germ
at every arity.

For covariance, the middle tails are related by
\[
 b_{(1)}(z)c_{(1)}\rho_{-z}(c_{(2)},b_{(2)})
       =\rho_{-z}(c_{(1)},b_{(1)})c_{(2)}b_{(2)}(z).
\]
This is \eqref{cylev:eq:two-positions}.
Multiplication on the left by $d_{(1)}$ and on the right by
$a_{(1)}(z)$ gives the desired reversed tensor coaction.

To prove skew-symmetry, expand the exchange pairing of
$b_{(1)}a_{(1)}$ with $d_{(1)}c_{(1)}$.
Alongside the two factor fields its scalar coefficient is
\[
\begin{split}
 &\rho_{-z}(c_{(1)},b_{(1)})\rho_z(b_{(2)},c_{(2)})
   \rho_z(b_{(3)},d_{(1)})\\
 &\hspace{18mm}\cdot\rho_z(a_{(1)},c_{(3)})
                         \rho_z(a_{(2)},d_{(2)}).
\end{split}
\]
The first two factors cancel. Skew-symmetry in each factor leaves
\[
 e^{z(T_1+T_2)}
 \bigl(Y_1(c,-z)a_{(0)}\otimes Y_2(d_{(0)},-z)b\bigr)
             \rho_z(a_{(1)},d_{(1)}),
\]
the reversed tensor field.
The proof of Theorem~\ref{cylev:thm:fields} now gives locality and
all exchange identities. Its pole bound uses the two factor bounds
and their finitely many cross pairings.

For associativity of the tensor itself, the reversed tails give
\begin{equation}\label{cylev:eq:tensor-annihilation}
 A_z^{V_1\boxtimes_\rho V_2}(x)
       =A_z^{V_1}(x_2)\otimes A_z^{V_2}(x_1).
\end{equation}
Both triple parenthesizations consequently have the field
\begin{equation}\label{cylev:eq:triple-field}
\begin{split}
 &Y_1(a,z)A_z^{V_1}(c_{(2)}b_{(1)})
     \otimes Y_2(b_{(0)},z)A_z^{V_2}(c_{(1)})\\
 &\hspace{47mm}\otimes Y_3(c_{(0)},z),
\end{split}
\end{equation}
and their coaction ends in $c_{(1)}b_{(1)}a_{(1)}$.
Thus the ordinary associator preserves the fields and coaction.
Its pentagon is the tensor-module pentagon.
A trivial coaction makes every cross factor a counit, proving the
unit assertion. Colinear field maps preserve the formulas.
\end{proof}

The tensor is also an instance of Li and Sun's twisted tensor
construction \cite[Definition~2.2 and Theorem~2.4]{LiSun2011}.
In their convention the twisting map is
\[
 \mathscr T(x)(b\otimes c)
       =c_{(0)}\otimes b_{(0)}\rho_x(c_{(1)},b_{(1)}).
\]
The two nontrivial twisting identities follow by expanding the tails
$d_{(1)}c_{(1)}(y)$ and $d_{(1)}b_{(1)}(y)$ in the covariance rule.
They read
\begin{align*}
 \mathscr T(x)(1\otimes Y_1(y))
   &=(Y_1(y)\otimes1)\mathscr T^{23}(x)\mathscr T^{12}(x+y),\\
 \mathscr T(x)(Y_2(y)\otimes1)
   &=(1\otimes Y_2(y))\mathscr T^{12}(x-y)\mathscr T^{23}(x).
\end{align*}
The unit identities are the two counits.
Their field formula $(Y_1(z)\otimes Y_2(z))\mathscr T^{23}(-z)$
is exactly \eqref{cylev:eq:tensor-field}.

\section{Ordered coproducts and covariant modules}

The tensor target has now been constructed independently of a level.
The coproduct into it is a shifted Hopf coproduct.

\begin{theorem}\label{cylev:thm:coproduct}
For any two levels, the map
\begin{equation}\label{cylev:eq:coproduct}
 D_{\eta_1,\eta_2}=(\eta_2\otimes1)\Delta:
 V_{\eta_1\eta_2}\longrightarrow V_{\eta_1}\boxtimes_\rho V_{\eta_2}
\end{equation}
preserves fields, vacuum, translation, coaction and exchange.
It is coassociative as a family:
\begin{equation}\label{cylev:eq:family-coassociativity}
 (D_{\eta_1,\eta_2}\boxtimes_\rho1)D_{\eta_1\eta_2,\eta_3}
 =(1\boxtimes_\rho D_{\eta_2,\eta_3})D_{\eta_1,\eta_2\eta_3}.
\end{equation}
The Hopf counit is a covariant field map
$\varepsilon:V_{\mathrm{id}}\to R$, and
\[
 (1\boxtimes_\rho\varepsilon)D_{\eta,\mathrm{id}}=1,
 \qquad
 (\varepsilon\boxtimes_\rho1)D_{\mathrm{id},\eta}=1.
\]
\end{theorem}
\begin{proof}
Let $D=(\eta_2\otimes1)\Delta$ and
$W=V_{\eta_1}\boxtimes_\rho V_{\eta_2}$.
Expanding its target coaction gives
\[
 \delta_WD(a)=\eta_2a_2\otimes a_5\otimes
 a_6\mathsf s^{-1}(\eta_2a_4)\eta_2a_3
                       \mathsf s^{-1}(\eta_1\eta_2a_1).
\]
The middle inverse-antipode pair cancels, leaving
$\eta_2a_2\otimes a_3\otimes
a_4\mathsf s^{-1}(\eta_1\eta_2a_1)$.
This is $(D\otimes1)\delta_{\eta_1\eta_2}(a)$.
The order of the levels has not changed.

Colinearity intertwines annihilations.
Using \eqref{cylev:eq:tensor-annihilation} and the multiplicativity
of $D$ for the original Hopf product gives
\begin{equation}\label{cylev:eq:coproduct-field}
 D Y_{\eta_1\eta_2}(a,z)
 =\bigl(L(\eta_2a_1(z))A_z^1(a_4)
           \otimes L(a_2(z))A_z^2(a_3)\bigr)D.
\end{equation}
Here $A^1,A^2$ denote the two target annihilations.

Expanding $Y_W(Da,z)$ gives in its first slot the factors
\[
 A_z^1(\eta_2a_2)A_z^1(a_6\mathsf s^{-1}(\eta_2a_3))
       =A_z^1(a_6\mathsf s^{-1}(\eta_2a_3)\eta_2a_2).
\]
Canceling the inverse-antipode pair gives the parenthesis in
\eqref{cylev:eq:coproduct-field}. Hence $D$ preserves fields.
It preserves the vacuum and translation by construction, and
colinearity also preserves the paired exchange.

The two iterates in \eqref{cylev:eq:family-coassociativity} both send
$a$ to
\begin{equation}\label{cylev:eq:triple-coproduct}
 \eta_2\eta_3a_1\otimes\eta_3a_2\otimes a_3.
\end{equation}
Theorem~\ref{cylev:thm:tensor} identifies their field targets.
At the identity level,
$(\varepsilon\otimes1)\delta_{\mathrm{id}}(b)=\varepsilon(b)1$.
Applying $\varepsilon$ to the field formula therefore gives
$\varepsilon(Y_{\mathrm{id}}(a,z)b)=\varepsilon(a)\varepsilon(b)$.
The same coaction identity preserves exchange.
The family counits follow from the Hopf counits and
$\varepsilon\eta=\varepsilon$.
\end{proof}

\begin{proposition}\label{cylev:prop:counit}
The Hopf counit $(H,\delta_\eta)\to R$, with trivial coaction on $R$,
is colinear if and only if $\eta=\mathrm{id}$.
Thus it is a covariant field counit precisely at the identity level.
\end{proposition}
\begin{proof}
Colinearity is the identity
$b_2\mathsf s^{-1}(\eta b_1)=\varepsilon(b)1$.
On $\operatorname{Hom}_R(H,H)$ use opposite convolution
$(f\star_{\mathrm{op}}g)(b)=f(b_2)g(b_1)$.
The displayed equation says that $\mathsf s^{-1}\eta$ is an inverse
of $\id$, whose inverse is already $\mathsf s^{-1}$.
Uniqueness gives $\eta=\id$.
The converse and preservation of fields are part of
Theorem~\ref{cylev:thm:coproduct}.
\end{proof}

Preservation of fields alone has a different answer.
Take $H=k[x]$ with $x$ primitive, $T=0$ and
$\rho_z=\varepsilon\otimes\varepsilon$.
The level $\eta x=2x$ gives ordinary multiplication as its field.
The polynomial augmentation preserves that field, but
$\delta_\eta(x)=x\otimes1-1\otimes x$, so it is not colinear.
Thus a scalar field character and a covariant counit are distinct.

For translation levels $\eta_i=e^{c_iT}$,
\[
 D_{c_1,c_2}(a)=a_1(c_2)\otimes a_2.
\]
Its source parameter is $c_1+c_2$, and its counit belongs to parameter
zero. The first target factor is translated by the second parameter.

\begin{definition}\label{cylev:def:module}
A \emph{covariant module} for an $H$-covariant field algebra $V$ is
a flat right $H$-comodule $M$ equipped with fields $Y_M(a,z)m\in M((z))$.
These satisfy the vacuum and input-translation identities.
Every finite ordered product on a fixed $m$ has a common germ
in $\mathcal G_r(M)$.
Every consecutive field substitution or terminal module substitution
expands that germ.

The terminal module is placed at $0$. Covariance means
\begin{equation}\label{cylev:eq:module-covariance}
 \delta_M(Y_M(a,z)m)
       =Y_M(a_{(0)},z)m_{(0)}\otimes m_{(1)}a_{(1)}(z).
\end{equation}
Morphisms preserve the fields and coaction.
Write $\mathcal M_\eta$ for these modules over $V_\eta$.
\end{definition}

No translation operator on the module itself is required here.
The identity $Y_M(Ta,z)=\partial_zY_M(a,z)$ concerns its state input.
Its coaction supplies the annihilation operators
\eqref{cylev:eq:annihilation} and the identity
\eqref{cylev:eq:annihilation-field}.

\begin{theorem}\label{cylev:thm:modules}
For covariant modules $M,N$ over $V_1,V_2$, respectively, the formula
\begin{equation}\label{cylev:eq:module-tensor}
\begin{split}
 &Y_{M\boxtimes_\rho N}(a\otimes b,z)(m\otimes n)\\
 &\qquad=Y_M(a,z)m_{(0)}\otimes Y_N(b_{(0)},z)n\,
                         \rho_{-z}(m_{(1)},b_{(1)})
\end{split}
\end{equation}
defines a covariant $V_1\boxtimes_\rho V_2$-module.
Its coaction has tail $n_{(1)}m_{(1)}$.
Pullback along \eqref{cylev:eq:coproduct} gives
\[
 \mathcal M_{\eta_1}\times\mathcal M_{\eta_2}
       \longrightarrow\mathcal M_{\eta_1\eta_2}.
\]
The disjoint union over levels is monoidal, with the ordinary
associativity map and unit $R\in\mathcal M_{\mathrm{id}}$ acting
through $\varepsilon$.
\end{theorem}
\begin{proof}
The action operator is
$Y_M(a,z)A_z^M(b_{(1)})\otimes Y_N(b_{(0)},z)$.
Moving annihilation past the next action produces
\eqref{cylev:eq:tensor-normal}, now acting on $M\otimes N$.
Module associativity in each factor identifies it with the action
of the tensor field. Repetition proves all consecutive substitutions.
Every pairing on fixed inputs is finite, so the same calculation
supplies the common denominator.
The unit rule gives the vacuum identity, and
$\partial_zA_z^M(x)=A_z^M(Tx)$ gives input translation.

Applying \eqref{cylev:eq:module-covariance} in both slots leaves the
same middle interchange \eqref{cylev:eq:two-positions} as in the
tensor-algebra proof. It gives the stated reversed tail.
For three modules both actions are \eqref{cylev:eq:triple-field}
with $Y_M,Y_N,Y_P$ in the respective slots.
Their pullbacks coincide by \eqref{cylev:eq:triple-coproduct}.

Consequently the usual associator is colinear and intertwines actions.
Its pentagon, the two family counits and the counital coactions give
all monoidal coherences. The formulas preserve module morphisms.
\end{proof}

The theorem uses the specified coaction on each module.
Without it, the annihilation $A_z^M$ in the tensor formula has no
definition. Likewise, the exchange on states does not by itself give
an intertwiner exchanging arbitrary modules.

\section{Polynomial fields and additive parameters}

Let $U$ be a free $R$-module and $H=\operatorname{Sym}_R(U)$, with
primitive generators and $T(U)\subset U$.
A bilinear kernel $G_z$ on $U$ extends uniquely by the unit and
multiplication rules to
\begin{equation}\label{cylev:eq:pairings}
 \rho_z(u_1\cdots u_r,v_1\cdots v_r)
       =\sum_{\sigma\in S_r}\prod_{i=1}^rG_z(u_i,v_{\sigma(i)}).
\end{equation}
The pairing of unequal polynomial degrees is zero.
Commutativity and cocommutativity give the form commutation rule.
Unitarity is equivalent to $G_z(u,v)+G_{-z}(v,u)=0$, and the
translation rule is its restriction to $U$.

For a level preserving $U$, put
\begin{equation}\label{cylev:eq:primitive-contraction}
 P_z^\eta(u,v)=\rho_{-z}(v-\eta v,u).
\end{equation}
Since $\delta_\eta(v)=v\otimes1+1\otimes(v-\eta v)$,
the primitive field is $Y_\eta(u,z)v=u(z)v+P_z^\eta(u,v)$.

For a basis $(u_i)$ and the corresponding polynomial derivatives,
all fields are given by Wick contraction:
\begin{equation}\label{cylev:eq:wick}
 Y_\eta(a,z)b=
 \mu(e^{zT}\otimes1)
 \exp\!\left(\sum_{i,j}P_z^\eta(u_i,u_j)
                         \partial_i\otimes\partial_j\right)(a\otimes b).
\end{equation}
Each input pair uses finitely many derivatives and finitely many
contractions. Selecting $r$ contracted factors in the coproducts and
pairing them gives \eqref{cylev:eq:pairings}. The exponential's
$1/r!$ cancels the ordering of the same contractions.
This proves the formula in every polynomial degree.

Take $H=R[x_1,x_2,\ldots]$, with $Tx_n=nx_{n+1}$.
An odd Laurent polynomial $g(z)$ defines the jet kernel
\begin{equation}\label{cylev:eq:jet-kernel}
 \rho_z(x_p,x_q)=\frac{(-1)^{q-1}}{(p-1)!(q-1)!}
                       \partial_z^{p+q-2}g(z).
\end{equation}
Indeed, multiplication by $p$ in the $(p+1,q)$ entry is its
$z$-derivative, and multiplication by $q$ in the $(p,q+1)$ entry
is its negative derivative.
Oddness gives the unitary condition.

For nilpotent $c$ and $\eta=e^{cT}$, the current $J=x_1$ has
\begin{align}
 P_c(z)&=g(z-c)-g(z),\label{cylev:eq:translation-kernel}\\
 \mathcal S_c(z)(J\otimes J)
 &=J\otimes J+
       \bigl(2g(z)-g(z+c)-g(z-c)\bigr)1\otimes1.
       \label{cylev:eq:translation-exchange}
\end{align}
The coproduct current is $J^{(1)}(c_2)+J^{(2)}$.
Its two self-contractions are $P_{c_1}$ and $P_{c_2}$, and its
mixed contraction is
\begin{equation}\label{cylev:eq:cross-kernel}
 Q_{c_1,c_2}(z)=g(z-c_1-c_2)-g(z-c_1)-g(z-c_2)+g(z).
\end{equation}
Their sum is $P_{c_1+c_2}$.
With $g(z)=\lambda/z$ and $c_1^2=c_2^2=0$, this is precisely
\eqref{cylev:eq:opening-cross}.

\subsection{An independent simple pole}

The double pole and simple pole of a current can vary independently.
Use $R=k[\kappa,\lambda]$ and two families of primitive jets
$x_n,y_n$, translated by $Tx_n=nx_{n+1}$ and $Ty_n=ny_{n+1}$.
At the first jets the kernel and level are
\begin{equation}\label{cylev:eq:two-current-data}
 G_z=\begin{pmatrix}
 -\lambda/(2z)&\kappa/z^2\\
 -\kappa/z^2&0
 \end{pmatrix},\qquad
 \eta x_n=-x_n+y_n,\quad \eta y_n=-y_n.
\end{equation}
Rows and columns of $G$ index its first and second inputs.
Extend each entry by the jet derivative formula and then by
\eqref{cylev:eq:pairings}.

\begin{proposition}\label{cylev:prop:two-current}
The data \eqref{cylev:eq:two-current-data} define a spectral form
and a level. The polynomial submodule generated by the $x_n$ is
a quantum vertex subalgebra, with primitive contraction
$\kappa z^{-2}+\lambda z^{-1}$ and exchange coefficient
$-2\lambda z^{-1}$ on $x_1\otimes x_1$.
\end{proposition}
\begin{proof}
The kernel satisfies $G_z+G_{-z}^{\mathsf t}=0$.
The matrix $E=\left(\begin{smallmatrix}-1&0\\1&-1\end{smallmatrix}\right)$
of the level satisfies $E^{\mathsf t}G_zE=G_z$ and commutes with
jet translation. Thus it preserves the extended form.

Since $x_1-\eta x_1=2x_1-y_1$,
\[
 \rho_{-z}(2x_1-y_1,x_1)=\frac\kappa{z^2}+\frac\lambda z,
 \qquad
 \rho_z(2x_1-y_1,2x_1-y_1)=-\frac{2\lambda}{z}.
\]
Differentiation gives all jet coefficients.
Wick contraction leaves only $x$-jets in the output, and the same
holds for exchange. This proves closure on all polynomial states.
\end{proof}

The full two-current algebra has the level-changing coproduct.
Its one-current subalgebra need not be preserved by the target levels:
already $\eta x_1=-x_1+y_1$ introduces the second current.
The larger state space is therefore part of this coproduct construction.
There is no vacuum-preserving $R$-linear field map from this singular
member to the scalar algebra. Such a map would send the nonzero
singular coefficients $\kappa1$ and $\lambda1$ to themselves,
whereas scalar fields have no pole.

\subsection{A noncommutative constant example}

The order of the tails can also be seen without a spectral parameter.
Over $k[\alpha]$, let $H$ have basis $1,g,x,gx$ with
\[
 g^2=1,\quad x^2=0,\quad xg=-gx,\qquad
 \Delta g=g\otimes g,\quad \Delta x=x\otimes1+g\otimes x.
\]
Its counit sends $g$ to $1$ and $x$ to $0$. Its antipode is
$\mathsf s(g)=g$, $\mathsf s(x)=-gx$.
Take $T=0$ and, in the displayed basis, the constant form
\begin{equation}\label{cylev:eq:sweedler-form}
 (\rho(e_i,e_j))=
 \begin{pmatrix}
 1&1&0&0\\1&-1&0&0\\0&0&\alpha&\alpha\\0&0&-\alpha&\alpha
 \end{pmatrix}.
\end{equation}
The multiplication and coproduct tables give both multiplicativity
rules, the commutation rule and unitarity.

The levels $\id$ and $\mathsf s^2$ have
\begin{align*}
 \delta_{\mathrm{id}}(x)&=x\otimes g+(1-g)\otimes gx,\\
 \delta_{\mathsf s^2}(x)&=x\otimes g-(1+g)\otimes gx.
\end{align*}
Substitution in \eqref{cylev:eq:field}, writing $*_{\eta}$ for the
constant field product, gives
\[
 x*_{\mathrm{id}}x=\alpha(1-g),\qquad
 x*_{\mathsf s^2}x=\alpha(1+g).
\]
The Hopf counit sends the second product to $2\alpha$.
Thus it fails to preserve fields at that level, while the identity-level
counit remains multiplicative.

\section{An integral formal family}

Let $\hbar$ be formal. A completed field on $V$ takes values in
\[
 V((z))^\wedge=\varprojlim_{N\geq1}(V/\hbar^NV)((z)).
\]
Each precision has a finite pole bound on fixed inputs.
The inverse limit need not have one bound valid at every precision.
All tensors in this section are $\hbar$-adically completed.

On $k[K][x_1,x_2,\ldots][[\hbar]]$, with the same jet translation,
consider the Wick propagator
\begin{equation}\label{cylev:eq:formal-propagator}
 P_K(z)=\frac1\hbar\left(\frac1{z-K\hbar}-\frac1z\right)
       =\frac K{z^2}+\frac{K^2\hbar}{z^3}
          +\frac{K^3\hbar^2}{z^4}+\cdots.
\end{equation}
Every fraction is expanded in $\hbar$.

\begin{proposition}\label{cylev:prop:formal}
These fields define compatible quantum vertex algebras at all finite
precisions, with Heisenberg classical limit of level $K$.
There are coassociative field maps from level $K_1+K_2$ given by
\begin{equation}\label{cylev:eq:formal-coproduct}
 x_n\longmapsto e^{K_2\hbar T}x_n\otimes1+1\otimes x_n.
\end{equation}
The target has diagonal propagators $P_{K_1},P_{K_2}$.
Its contraction from the first current to the second is zero.
The contraction from the second to the first is
\begin{equation}\label{cylev:eq:formal-cross}
 Q_{K_1,K_2}^{21}(z)=\frac1\hbar
 \left(-\frac1z+\frac1{z+K_2\hbar}+\frac1{z-K_1\hbar}
                 -\frac1{z+(K_2-K_1)\hbar}\right).
\end{equation}
The family counit is polynomial augmentation at $K=0$.
\end{proposition}
\begin{proof}
The numerator in \eqref{cylev:eq:formal-propagator} has zero constant
term. In \eqref{cylev:eq:formal-cross}, both constant and linear terms
cancel, so $Q^{21}=O(\hbar)$.
The displayed fields are therefore integral.

At finite precision every propagator is a Laurent polynomial.
The contraction operators in the Wick formula commute.
A product and its iterate contract exactly the same pairs, giving
weak associativity on a common germ.

Exchanging the two inputs replaces the contraction matrix by its
transposed value at $-z$. The exponential of the difference gives
exchange, with unitarity, Yang--Baxter and the hexagon following
from addition of these commuting pair contractions.
Jet differentiation gives translation.
Reduction commutes with all these finite formulas.
At $\hbar=0$ the propagator is $K/z^2$ and the exchange is the identity.

The contraction of the image current is
\[
 P_{K_1}(z)+P_{K_2}(z)+Q_{K_1,K_2}^{21}(z-K_2\hbar)
       =P_{K_1+K_2}(z).
\]
The four fractions cancel to give this identity.
The state map preserves translation and polynomial multiplication,
so Wick contraction proves preservation of every field and exchange.

Writing $Q_{A,B}=Q_{A,B}^{21}$, the same fractions give
\begin{align*}
 Q_{K_1+K_2,K_3}(z)
   &=Q_{K_1,K_3}(z-K_2\hbar)+Q_{K_2,K_3}(z),\\
 Q_{K_1,K_2+K_3}(z)
   &=Q_{K_1,K_2}(z+K_3\hbar)+Q_{K_1,K_3}(z).
\end{align*}
Thus both triple parenthesizations have the same pairwise kernels.
Both iterated state maps shift the first factor by
$(K_2+K_3)\hbar$, the second by $K_3\hbar$, and the third by zero.
This proves coassociativity at every precision and in the limit.
At $K=0$ the fields are regular polynomial fields, whose augmentation
gives the counit.
\end{proof}

Formally, the primitive spectral form behind this example is
$1/(\hbar z)$, which is not integral.
The difference between a state and its level translate supplies the
power of $\hbar$ canceled in the field formula.
Thus integrality of fields and integrality of the underlying form
are separate assertions.

\section{Divided states and the rational quantum affine algebra}

The same distinction occurs for matrix currents.
Write $\mathsf R$ for an auxiliary matrix, reserving $R$ for the
coefficient ring. If $L=1+\hbar\ell$ and
$\rho_z(L_1(u),L_2(v))=\mathsf R(z+u-v)$, the unit rule forces
\begin{equation}\label{cylev:eq:divided-form}
 \rho_z(\ell_1(u),\ell_2(v))
       =\hbar^{-2}\bigl(\mathsf R(z+u-v)-1\bigr).
\end{equation}
Thus $\mathsf R=1+\hbar r+O(\hbar^2)$ with $r\ne0$ gives a
term $\hbar^{-1}r$ on divided generators.
The formula follows simply by subtracting the identity in both inputs.
An integral field construction must explain how its actual coaction
avoids this nonintegral pairing.

Work now over $\mathbb C[[\hbar]]$ in rank two.
Let $P$ flip $\mathbb C^2\otimes\mathbb C^2$, and normalize
\begin{equation}\label{cylev:eq:rtt-normalization}
\begin{split}
 \mathsf R(w)&=f(\hbar/w)(1-\hbar P/w),\qquad f(0)=1,\\
 f\left(\frac{\hbar}{w+2\hbar}\right)
       &=f(\hbar/w)(1-\hbar^2/w^2).
\end{split}
\end{equation}
This determines the scalar series, whose first terms are
\[
 f(t)=1+\frac12t+\frac58t^2+\frac{11}{16}t^3
                         +\frac{83}{128}t^4+O(t^5).
\]
The matrix is unitary and has crossing identity
\begin{equation}\label{cylev:eq:rtt-crossing}
 \bigl((\mathsf R(w)^{t_1})^{-1}\bigr)^{t_1}
       =\mathsf R(w-2\hbar)^{-1},
\end{equation}
where $t_1$ transposes the first auxiliary factor.
The scalar equation and $(P^{t_1})^2=2P^{t_1}$ prove these identities.

Let $H$ be the $\hbar$-adic dual Yangian generated by
\[
 L(u)=1+\hbar\ell(u),\qquad
 \ell(u)=\sum_{r\geq1}\ell^{(r)}u^{r-1},
\]
with RTT relation
\begin{equation}\label{cylev:eq:rtt}
 \mathsf R(u-v)L_1(u)L_2(v)
       =L_2(v)L_1(u)\mathsf R(u-v).
\end{equation}
Its generator relations are obtained after canceling $\hbar^2$.
We use the torsion-free algebra with its topological PBW basis of
ordered monomials in the $\ell_{ij}^{(r)}$.

The coproduct is matrix multiplication, $\Delta L=L^{[1]}L^{[2]}$,
the counit sends $L$ to $1$, and $TL(u)=\partial_uL(u)$.
The internal variables $u,v$ are expanded at zero.
The normalized matrix defines the spectral form on polynomial states
after inverting $\hbar$. The multiplication rules determine every
pairing.

The independent comparison object is the double Yangian at level $c$.
It has a creation matrix
$T^+(u)=1-\hbar\sum_{r\geq1}t^{(-r)}u^{r-1}$ and an annihilation
matrix $M_c(s)=1+\hbar\sum_{r\geq1}m^{(r)}s^{-r}$.
Each satisfies RTT, and their mixed relation is
\begin{equation}\label{cylev:eq:double-mixed}
 \mathsf R(s-v+c\hbar/2)M_{c,1}(s)T^+_2(v)
       =T^+_2(v)M_{c,1}(s)\mathsf R(s-v-c\hbar/2).
\end{equation}
It is expanded in $s^{-1}$ and $v$ and divided by $\hbar^2$
at the level of generator relations.

The vacuum module $\mathcal V_c$ is the completed quotient by the
left ideal generated by the $m_{ij}^{(r)}$, so
$M_c(s)\mathbf1=1$.
The PBW theorem identifies its states with the completed dual Yangian
\cite[Theorem~2.2 and \S4.1]{JingKozicMolevYang2017}.
The full scalar-matrix direction is retained.

\begin{lemma}\label{cylev:lem:integrality}
For $\eta_c=e^{-c\hbar T}$, the fields and exchange of $V_{\eta_c}$
are integral on the divided states $H$.
The cross contractions in
$V_{\eta_{c_1}}\boxtimes_\rho V_{\eta_{c_2}}$ are integral as well.
They extend continuously to the $\hbar$-adic state completions.
\end{lemma}
\begin{proof}
Let $F_pH$ be PBW degree at most $p$ before completion.
The completed subalgebra $B$ generated by entries of $L$ consists,
in PBW coordinates, of series
$\sum_{r\geq0}\hbar^rb_r$ with $b_r\in F_rH$.
The divided RTT relations give
$[\ell_{ij}^{(r)},\ell_{kl}^{(s)}]\in F_1H+\hbar F_2H$.
Consequently $[\ell_{ij}^{(r)},B]\subset B$.
Both antipodes, the coproduct and translation preserve $B$, and
$(\eta_c-1)B\subset\hbar B$.

The pairing has the estimate
\begin{equation}\label{cylev:eq:degree-estimate}
 \rho_z(F_pH,F_qH)
       \subset\hbar^{-\min(p,q)}\mathbb C((z))[[\hbar]].
\end{equation}
To see it, pair two products of $L$'s and subtract the identity
in every input. Each surviving product of factors $\mathsf R-1$
must meet every input generator. At least $\max(p,q)$ such factors
are needed. Dividing by $\hbar^{p+q}$ gives the estimate.

It makes pairings of $B$ with polynomial states integral.
For an input of degree $q$, the term $\hbar^rb_r$ contributes
at order at least $r-q$, proving convergence.

The coaction has its second component in $B$.
On a divided generator the inverse-antipode identity gives
\begin{equation}\label{cylev:eq:integral-coaction}
\begin{split}
 \delta_{\eta_c}(\ell_{ij}(u))
 ={}&\sum_{a,b}\ell_{ab}(u)\otimes
                  L_{bj}(u)\mathsf s^{-1}\eta_cL_{ia}(u)\\
 &+1\otimes\hbar^{-1}\sum_aL_{aj}(u)
                  \mathsf s^{-1}(\eta_c-1)L_{ia}(u).
\end{split}
\end{equation}
Both tails lie in $B$.
For products, \eqref{cylev:eq:coaction-product} introduces only
commutators $[\ell_{ij}^{(r)},b]$ and the divided translation
difference in the second line, multiplied by elements of $B$.
Induction gives $\delta_{\eta_c}(H)\subset H\widehat\otimes B$.

Annihilation therefore pairs $B$ with a polynomial state.
Exchange and cross contraction pair two tails in $B$.
The estimate proves integrality and convergence for each operation.
Its continuous $\hbar$-linear extension gives the completed fields.
At finite precision a fixed tuple of states still has a finite
Laurent pole bound.
\end{proof}

\begin{theorem}\label{cylev:thm:rtt-comparison}
The map
\begin{equation}\label{cylev:eq:rtt-state-map}
 \Phi_c(\ell_{i_1j_1}^{(r_1)}\cdots\ell_{i_mj_m}^{(r_m)})
       =(-1)^m t_{i_1j_1}^{(-r_1)}\cdots
                        t_{i_mj_m}^{(-r_m)}\mathbf1
\end{equation}
identifies $V_{e^{-c\hbar T}}$ with the rational quantum affine
vertex algebra on $\mathcal V_c$.
It preserves the vacuum, translation, all fields and exchange.
On products of generating matrices the fields are
\begin{equation}\label{cylev:eq:rtt-fields}
\begin{split}
 Y_c(L_1(u_1)\cdots L_m(u_m)\mathbf1,z)
 ={}&L_1(z+u_1)\cdots L_m(z+u_m)\\
 &\cdot\bigl(M_{c,1}(z+u_1+c\hbar/2)\cdots
               M_{c,m}(z+u_m+c\hbar/2)\bigr)^{-1}.
\end{split}
\end{equation}
\end{theorem}
\begin{proof}
PBW gives the continuous inverse replacing $t_{ij}^{(-r)}$ by
$-\ell_{ij}^{(r)}$.
Differentiation of generating series proves translation compatibility.
Put $K_0(u,z)=A_z(L_0(u))$ and $w=z+u-v$.
The antipode pairing rule and crossing identity give
$\rho_w(\mathsf s^{-1}L_1,L_2)=\mathsf R(w-2\hbar)^{-1}$.
Substitution in \eqref{cylev:eq:mixed-operators} yields
\begin{equation}\label{cylev:eq:annihilation-recurrence}
 K_0(u,z)\mathsf R_{01}(w+c\hbar)^{-1}L_1(v)
       =L_1(v)\mathsf R_{01}(w)^{-1}K_0(u,z),
 \qquad K_0(u,z)\mathbf1=1.
\end{equation}
Inverting \eqref{cylev:eq:double-mixed} shows that
$M_c(z+u+c\hbar/2)^{-1}$ satisfies this same recurrence.

Its value on the vacuum determines its action on every ordered
product of creation matrices. Subtracting their constant terms and
dividing by powers of $\hbar$ gives equality on polynomial PBW states.
Lemma~\ref{cylev:lem:integrality} extends it to the completion.
The anti-action reverses annihilation order, giving
\eqref{cylev:eq:rtt-fields}.

The exchange uses four contractions of the two coaction tails.
For an auxiliary matrix $F=\sum F'\otimes F''$, write
$F^{\mathrm{lr}}(X\otimes Y)=\sum F'X\otimes YF''$.
The superscripts $\mathrm{ll},\mathrm{rl},\mathrm{rr}$ indicate the
other left or right placements.

Expanding \eqref{cylev:eq:exchange}, using the two multiplicativity
rules and the inverse-antipode pairing, gives the operator
\begin{equation}\label{cylev:eq:four-contractions}
 \bigl(\mathsf R(w-c\hbar)^{\mathrm{lr}}\bigr)^{-1}
       \mathsf R(w)^{\mathrm{ll}}\mathsf R(w)^{\mathrm{rr}}
       \bigl(\mathsf R(w+c\hbar)^{-1}\bigr)^{\mathrm{rl}}.
\end{equation}
The first inverse is an inverse of an operator.
The last factor places an inverse matrix. These two operations differ.

Equivalently,
\begin{equation}\label{cylev:eq:exchange-comparison}
\begin{split}
 &\mathcal S_c(z)\bigl(\mathsf R(w)^{-1}L_2(v)
                       \mathsf R(w-c\hbar)L_1(u)\bigr)\\
 &\qquad=L_1(u)\mathsf R(w+c\hbar)^{-1}L_2(v)\mathsf R(w).
\end{split}
\end{equation}
The state factors remain separate throughout the auxiliary products.
This is the exchange on the independently defined vacuum module
\cite[Theorem~4.1, equations~(4.17)--(4.19)]{JingKozicMolevYang2017}.

For $n$ matrices in the first state and $m$ in the second, replace
each $\mathsf R(w)$ by
$\prod_{i=1}^{n}\prod_{j=m}^{1}
\mathsf R_{i,n+j}(z+u_i-v_j)$, in those orders.
Appending one creation matrix inserts its four contractions by
the multiplicativity rules. Yang--Baxter moves them past the
existing factors, and crossing cancels the inverse pairs.
Induction proves the same exchange identity for all polynomial
states. Continuity proves it for completed states.
\end{proof}

The first divided coefficients make the level visible.
Put $t=\hbar/w$ and let
$A=P^{\mathrm{lr}}$, $B=P^{\mathrm{ll}}$,
$C=P^{\mathrm{rr}}$, $D=P^{\mathrm{rl}}$.
Then $A^2=2A$, $D^2=2D$ and $B^2=C^2=1$.
Writing $a=c+2$, expansion of \eqref{cylev:eq:four-contractions} is
\begin{equation}\label{cylev:eq:fourth-order}
\begin{split}
 \mathcal Q_c(t)={}&(1+t^2-c(c+2)t^3+t^4)\\
 &\cdot(1+(t+at^2+a^2t^3+a^3t^4)A)\\
 &\cdot(1-t(B+C)+t^2BC)\\
 &\cdot(1+(t-ct^2+c^2t^3-c^3t^4)D)\pmod{t^5}.
\end{split}
\end{equation}
The scalar factor comes from
\[
 \frac{f(t)^2}
 {f(t/(1-ct))f(t/(1+ct))(1-t^2/(1+ct)^2)}.
\]

Let $Q_j$ be the coefficient of $t^j$.
For $X=\ell_1(u)+\ell_2(v)$ and $Z=\ell_1(u)\ell_2(v)$,
subtracting the vacuum and one-state terms of $L_1L_2$ gives
\begin{equation}\label{cylev:eq:divided-exchange}
\begin{split}
 \mathcal S_c(z)Z={}&Z+
 \hbar\left(\frac{Q_1Z}{w}+\frac{Q_2X}{w^2}+\frac{Q_3(1)}{w^3}\right)\\
 &+\hbar^2\left(\frac{Q_2Z}{w^2}+\frac{Q_3X}{w^3}
                                      +\frac{Q_4(1)}{w^4}\right)
       +O(\hbar^3).
\end{split}
\end{equation}
Direct multiplication gives $Q_1(1)=Q_2(1)=Q_1X=0$ and, with
$\Omega=P-\tfrac12I$,
\begin{equation}\label{cylev:eq:vacuum-coefficients}
 Q_3(1)=2c^2\Omega,\qquad Q_4(1)=4c^2\Omega.
\end{equation}
For the $11,11$ input, applying both Hopf counits gives
$c^2\hbar w^{-3}+2c^2\hbar^2w^{-4}+O(\hbar^3)$.
Replacing the level translate by the identity sets $c=0$ and loses
this coefficient already at first order on divided states.

The fields retain the classical affine level as well.
For $J_{ij}(z)=Y_c(-\ell_{ij}^{(1)},z)$, expansion of the mixed
relation through $\hbar^2$ gives, modulo $\hbar$,
\begin{equation}\label{cylev:eq:classical-affine}
 J_{ij}(z)J_{kl}(w)\sim
 \frac{c(\delta_{jk}\delta_{il}-\tfrac12\delta_{ij}\delta_{kl})}{(z-w)^2}
 +\frac{\delta_{jk}J_{il}(w)-\delta_{il}J_{kj}(w)}{z-w}.
\end{equation}
The difference of the two shifted $R$-matrices supplies the double
pole, while their linear matrix terms supply the Lie bracket.
The invariant form is
$\operatorname{tr}(XY)-\tfrac12\operatorname{tr}(X)\operatorname{tr}(Y)$.
Its scalar-matrix direction has zero level.
For $c\ne0$, the $12,21$ double pole excludes a vacuum-preserving
field map to the scalar vertex algebra.

\section{The two rational coproducts}

Under \eqref{cylev:eq:rtt-state-map}, the spectral coproduct is
\begin{align}
 \mathscr D_{c_1,c_2}L(u)
 &=L^{[1]}(u-c_2\hbar)L^{[2]}(u),
       \label{cylev:eq:spectral-coproduct}\\
 \mathscr D_{c_1,c_2}\ell(u)
 &=\ell^{[1]}(u-c_2\hbar)+\ell^{[2]}(u)
       +\hbar\ell^{[1]}(u-c_2\hbar)\ell^{[2]}(u).
       \nonumber
\end{align}
Its source is $\mathcal V_{c_1+c_2}$ and its target is
$\mathcal V_{c_1}\boxtimes_\rho\mathcal V_{c_2}$.
Theorems~\ref{cylev:thm:coproduct} and
\ref{cylev:thm:rtt-comparison} prove that it preserves all fields
and exchange. Both triple iterates give
\[
 L^{[1]}(u-(c_2+c_3)\hbar)L^{[2]}(u-c_3\hbar)L^{[3]}(u).
\]
The counit belongs to level zero.

The field of its generating state is
\begin{equation}\label{cylev:eq:coproduct-fields}
\begin{split}
 Y_{12}(\mathscr D L(u),z)
 ={}&L^{[1]}(u+z-c_2\hbar)L^{[2]}(u+z)\\
 &\cdot M_{c_2}^{[2]}(u+z+c_2\hbar/2)^{-1}
          M_{c_1}^{[1]}(u+z+c_1\hbar/2)^{-1}.
\end{split}
\end{equation}
This follows by inserting the annihilation recurrence into
\eqref{cylev:eq:coproduct-field}.
The first creation matrix is shifted by $-c_2\hbar$, while its
annihilation argument has no such shift. The cross contraction
accounts for the difference.

The ordinary double Yangian has a different pair of half-coproducts.
With central generator $C$, put $C_1=C\otimes1$, $C_2=1\otimes C$.
In the normalization of \eqref{cylev:eq:double-mixed}, they are
\begin{align}
 \Delta_{\mathrm{DY}}L(u)
 &=L^{[1]}(u-\hbar C_2/4)L^{[2]}(u+\hbar C_1/4),
       \label{cylev:eq:dy-creation}\\
 \Delta_{\mathrm{DY}}M(s)
 &=M^{[1]}(s+\hbar C_2/4)M^{[2]}(s-\hbar C_1/4),
       \label{cylev:eq:dy-annihilation}\\
 \Delta_{\mathrm{DY}}C&=C_1+C_2.\nonumber
\end{align}
These opposite quarter shifts are part of the Hopf structure
\cite[\S4.1]{JingKozicMolevYang2017}.

Scalar recentering of each level cannot identify the families.
If $L(u)$ is replaced by $L(u+b(c)\hbar)$ at level $c$, the two
spectral shifts become
$b(c_1)-b(c_1+c_2)-c_2$ and $b(c_2)-b(c_1+c_2)$.
Equality with the quarter shifts, first at $c_1=0$ and then at
$c_2=0$, would require simultaneously
\[
 b(c)-b(0)=-3c/4,\qquad b(c)-b(0)=-c/4.
\]
A tensor-dependent translation does give an exact comparison.

\begin{definition}\label{cylev:def:translation-extension}
Adjoin to the double Yangian an element $d$ with
\[
 [d,C]=0,\qquad[d,L(u)]=\partial_uL(u),\qquad
 [d,M(s)]=\partial_sM(s),
\]
and make $d$ primitive, with counit zero and antipode $-d$.
This is the \emph{translation extension} of the double Yangian.
\end{definition}

Simultaneous differentiation leaves every spectral difference in the
defining relations unchanged, so it defines a derivation.
Differentiating the two half-coproducts proves the coderivation rule.
The extension therefore exists as a completed Hopf algebra.
On $\mathcal V_c$, $C$ acts by $c$ and $d$ by the state translation $T$.

\begin{theorem}\label{cylev:thm:twist}
In the translation extension,
\begin{equation}\label{cylev:eq:twist}
 F=\exp\!\left(\frac{\hbar}{4}(3d\otimes C+C\otimes d)\right)
\end{equation}
is an invertible normalized Hopf two-cocycle.
The conjugated coproduct $\Delta_{\mathrm{sp}}=F^{-1}\Delta_{\mathrm{DY}}F$
satisfies
\begin{align}
 \Delta_{\mathrm{sp}}L(u)&=L^{[1]}(u-\hbar C_2)L^{[2]}(u),
       \label{cylev:eq:sp-creation}\\
 \Delta_{\mathrm{sp}}M(s)&=M^{[1]}(s-\hbar C_2/2)
                              M^{[2]}(s-\hbar C_1/2).
       \label{cylev:eq:sp-annihilation}
\end{align}
It restricts to the unextended double.
On vacuum state spaces,
\begin{equation}\label{cylev:eq:state-twist}
 F_{c_1,c_2}=e^{\hbar(3c_2T_1+c_1T_2)/4},\qquad
 \Delta^{\mathrm{vac}}_{\mathrm{DY},c_1,c_2}
       =F_{c_1,c_2}\mathscr D_{c_1,c_2}.
\end{equation}
This is a continuous invertible comparison preserving the vacuum
and total translation.
\end{theorem}
\begin{proof}
A normalized two-cocycle satisfies
\[
 F_{12}(\Delta\otimes1)F=F_{23}(1\otimes\Delta)F,
 \qquad(\varepsilon\otimes1)F=(1\otimes\varepsilon)F=1.
\]
Here $C,d$ commute and are primitive.
For $f=(3d\otimes C+C\otimes d)/4$, all placements commute and
$f_{12}+(\Delta\otimes1)f=f_{23}+(1\otimes\Delta)f$.
Exponentiation proves the two-cocycle equation and counitality.
The same argument applies to $F^{-1}$.

Conjugation by $F^{-1}$ shifts first-factor arguments by
$-3\hbar C_2/4$ and second-factor arguments by $-\hbar C_1/4$.
Applied to the two opposite quarter shifts, it gives
\eqref{cylev:eq:sp-creation}--\eqref{cylev:eq:sp-annihilation}.
Their coefficients lie in the unextended tensor square.
The twisted antipode is conjugated by
$m(1\otimes\mathsf s_{\mathrm{DY}})(F^{-1})=e^{\hbar Cd}$.
Its action translates both matrix arguments by $\hbar C$, so it
also preserves the unextended algebra.

Both half-coproducts make the annihilation matrices act as the identity
on $\mathbf1\otimes\mathbf1$ and therefore descend to the vacuum
quotients of source level $c_1+c_2$.
There $F$ acts as the exponential in \eqref{cylev:eq:state-twist}.
Conjugation proves the state identity on every creation generator,
then on their products and their PBW completion.
\end{proof}

The comparison transports the field tensor as well.
Put $p=3\hbar c_2/4$, $q=\hbar c_1/4$ and $\tau_a=e^{aT}$.
On the same completed tensor module, set
\begin{align}
 Y^F(a\otimes b,z)(x\otimes y)
 &=Y_{c_1}(a,z)x_{(0)}\otimes Y_{c_2}(b_{(0)},z)y\,
       \rho_{-z-p+q}(x_{(1)},b_{(1)}),\label{cylev:eq:twisted-field}\\
 \delta^F(a\otimes b)
 &=a_{(0)}\otimes b_{(0)}\otimes
          \tau_{-q}(b_{(1)})\tau_{-p}(a_{(1)}).
       \label{cylev:eq:twisted-coaction}
\end{align}
The shifts are $\hbar$-adic Taylor series.
The pairings use the integral tails of Lemma~\ref{cylev:lem:integrality}.

\begin{proposition}\label{cylev:prop:transported-tensor}
The map $F_{c_1,c_2}=\tau_p\otimes\tau_q$ is an isomorphism
from $\mathcal V_{c_1}\boxtimes_\rho\mathcal V_{c_2}$ to the
field tensor \eqref{cylev:eq:twisted-field}--\eqref{cylev:eq:twisted-coaction}.
It preserves fields, vacuum, translation, coaction and exchange.
In this target, the ordinary double-Yangian vacuum coproduct is a
field homomorphism. The comparison is coherent on all iterated tensors.
\end{proposition}
\begin{proof}
Translation is a field automorphism, and
$\delta\tau_a=(\tau_a\otimes\tau_a)\delta$.
The form satisfies
$\rho_w(\tau_a x,\tau_b y)=\rho_{w+a-b}(x,y)$.
Inserting the translated inputs in \eqref{cylev:eq:twisted-field}
therefore changes its scalar factor to
$\rho_{-z-p+q+p-q}=\rho_{-z}$.
All other factors are the translates of the original fields.
Thus $Y^F(Fv,z)Fw=FY_{12}(v,z)w$.
The same substitution gives $\delta^FF=(F\otimes1)\delta_{12}$,
and hence preservation of the paired exchange.

These identities hold at every precision, so also on the completion.

Composing with $\mathscr D$ and using \eqref{cylev:eq:state-twist}
gives the ordinary vacuum coproduct as a field map.
On three factors both comparison composites are
\[
 \exp\!\left(\frac{\hbar}{4}
  \bigl(3(c_2+c_3)T_1+(c_1+3c_3)T_2+(c_1+c_2)T_3\bigr)\right).
\]
This is the cocycle identity on state spaces.
It gives the associativity coherence, and iteration gives all higher
coherences. The scalar unit has level and translation zero, so both
unit comparisons are the identity.
\end{proof}

\subsection{Obstructions without translation}

The translation extension changes the admissible comparison maps.
This can be detected by a central series in the unextended double.
Define the quantum determinant
\begin{equation}\label{cylev:eq:determinant}
 Q(u)=L_{11}(u)L_{22}(u-\hbar)-L_{21}(u)L_{12}(u-\hbar).
\end{equation}
Its coefficients are central
\cite[Proposition~2.8]{JingKozicMolevYang2017}.
For commuting RTT matrices it is multiplicative: the identity
$A^{(2)}L_1(u)L_2(u-\hbar)=A^{(2)}Q(u)$, with
$A^{(2)}=(1-P)/2$, applied successively to the two matrices proves it.

\begin{proposition}\label{cylev:prop:central-obstruction}
At a nonzero level pair $(c_1,c_2)$, no invertible element of the
unextended double-Yangian tensor square conjugates the spectral
coproduct to the ordinary double-Yangian coproduct.
The obstruction is already visible modulo $\hbar^3$.
\end{proposition}
\begin{proof}
The images of the quantum determinant are
\begin{align*}
 \Delta_{\mathrm{sp}}Q(u)&=Q^{[1]}(u-c_2\hbar)Q^{[2]}(u),\\
 \Delta_{\mathrm{DY}}Q(u)&=Q^{[1]}(u-c_2\hbar/4)
                                  Q^{[2]}(u+c_1\hbar/4).
\end{align*}
Both are central in the target and hence fixed by conjugation.
Since $L=1+\hbar\ell$,
$\Delta_{\mathrm{DY}}Q-\Delta_{\mathrm{sp}}Q$ begins with
\begin{equation}\label{cylev:eq:central-defect}
 \frac{\hbar^2}{4}
 \left(3c_2\partial_u\operatorname{tr}\ell^{[1]}(u)
          +c_1\partial_u\operatorname{tr}\ell^{[2]}(u)\right).
\end{equation}
PBW makes the two trace derivatives independent, so this coefficient
is nonzero for every nonzero level pair.

For the scalar test $c_1=0,c_2=1$, take the creation-algebra characters
$L^{[1]}(u)=(1+\hbar u)I_2$, $L^{[2]}(u)=I_2$.
The coefficient in \eqref{cylev:eq:central-defect} is $3/2$.
Only the creation subalgebras are evaluated in this test. The
characters need not extend to the annihilation matrices.
\end{proof}

At $(c_1,c_2)=(0,0)$ the state coproducts agree.
The determinant-one quotient removes this particular central coefficient.
In the translation extension, $[d,Q(u)]=\partial_uQ(u)$,
so this coefficient is no longer central and the explicit twist applies.

A different restriction concerns the completion.
Let
\[
 A_c=U(\mathfrak{sl}_2[t,t^{-1}]\oplus\mathbb CK)/(K-c),
\]
where
$[X_m,Y_n]=[X,Y]_{m+n}
+m\delta_{m,-n}\operatorname{tr}(XY)K$.
An element of $A_c$ is a finite sum of PBW monomials.
No completion in the mode index is included.
Take $H=\operatorname{diag}(1,-1)$, $E=E_{12}$ and $F=E_{21}$.

\begin{lemma}\label{cylev:lem:finite-mode}
For $(\alpha,\beta)\ne(0,0)$, no
$a\in A_{c_1}\otimes A_{c_2}$ satisfies, for every $r\geq1$,
\[
 [a,H_{-r}\otimes1+1\otimes H_{-r}]
       =r(\alpha H_{-r-1}\otimes1+\beta1\otimes H_{-r-1}).
\]
\end{lemma}
\begin{proof}
Use PBW symmetrization to express $a$ as a finite polynomial in
loop generators in each slot.
Choose $r$ larger than the absolute value of every mode index in it.
Then its commutator with $H_{-r}$ has no central term.
The adjoint action sends an $E_n$ or $F_n$ to a root vector of mode
$n-r$ and sends $H_n$ to zero.
Symmetrization intertwines this action, so it preserves PBW bidegree.

Only bidegrees $(1,0)$ and $(0,1)$ could produce the right side.
Their brackets are root vectors, while that side is a nonzero
combination of independent Cartan vectors. This is impossible.
\end{proof}

Consequently, suppose a flat determinant-one quantization has
classical reduction $A_c$ and translation
$[d,H_{-r}]=rH_{-r-1}$.
A twist $1+\hbar a+O(\hbar^2)$ in its purely $\hbar$-adic tensor
completion would give the forbidden identity with
$\alpha=3c_2/4$, $\beta=c_1/4$.
This consequence assumes the specified classical reduction.
A completion allowing infinitely many modes in $a$ has a different
admissible class.

Such mode completions occur in the double-Yangian cocycle of
Enriquez and Felder \cite[\S1, Theorem~3.1 and Proposition~5.1]{EnriquezFelder1997}.
Their algebra includes a derivation and central element, and their
cocycle compares Drinfeld's coproduct with the Yangian coproduct.
Those two named structures require their own current and completion
comparison with \eqref{cylev:eq:dy-creation}--\eqref{cylev:eq:sp-annihilation}.
The finite-mode lemma does not exclude their construction.

\subsection{The field tensor remains nontrivial}

Even when the two state coproducts agree at level zero, their field
target retains a mixed exchange.
Let $j_1(a)=a\otimes\mathbf1$, $j_2(b)=\mathbf1\otimes b$ in
$W=\mathcal V_0\boxtimes_\rho\mathcal V_0$.
Write its exchange as $1+\hbar\sigma_W(z)+O(\hbar^2)$.
Use the ordered PBW decomposition to project to degree one in each
state factor, denoting this projection by $\operatorname{pr}_{1,1}$.

\begin{proposition}\label{cylev:prop:ordinary-obstruction}
There is no exchange-preserving isomorphism
$W\to\mathcal V_0\otimes\mathcal V_0$ equal to the identity
modulo $\hbar$, when the target has ordinary tensor fields and exchange.
\end{proposition}
\begin{proof}
The four first-order contractions in the mixed exchange give
\begin{equation}\label{cylev:eq:mixed-witness}
 \operatorname{pr}_{1,1}\sigma_W(z)
       (\ell_{11}^{(1)}\otimes\ell_{12}^{(1)})
       =z^{-1}\ell_{12}^{(1)}\otimes
                          (\ell_{22}^{(1)}-\ell_{11}^{(1)}).
\end{equation}
The inputs denote $j_1(\ell_{11}^{(1)})$ and
$j_2(\ell_{12}^{(1)})$ in the two copies of $W$.
In auxiliary notation their coefficient is
$P^{\mathrm{lr}}Z-PZ-ZP+P^{\mathrm{rl}}Z$,
where $Z=\ell_1\ell_2$.
Its $11,12$ entry is the displayed tensor.

Equivalently, the adjoint Casimir calculation is
\[
 -\sum_{i,j}[E_{ij},E_{11}]\otimes[E_{ji},E_{12}]
       =E_{12}\otimes(E_{22}-E_{11}).
\]
PBW makes \eqref{cylev:eq:mixed-witness} nonzero.
The ordinary tensor exchange is the identity on these mixed inputs.

If $G=1+\hbar g+O(\hbar^2)$ intertwined the exchanges, the coefficient
of $\hbar$ in
$\mathcal S_{\mathrm{ord}}(G\otimes G)=(G\otimes G)\mathcal S_W$
would give $\sigma_{\mathrm{ord}}=\sigma_W$.
The nonzero mixed coefficient contradicts this equality.
\end{proof}

Thus the translation twist compares two presentations of the same
level family with their specified field tensors.
It does not identify that tensor with the ordinary field tensor.
The obstruction just proved fixes the identity classical state map.
It makes no assertion about every possible intrinsic comparison.

\section{Differentials and the spectral tensor of a line}
\label{cylev:sec:dg}

A spectral tensor can also change a differential.
For a dg algebra $(A,Q)$, an element $r(z)$ of degree one with
$Q_\otimes r+r^2=0$ defines the differential
$Q_\otimes+[r(z),-]$ on the completed tensor square containing $r$.
Here $Q_\otimes=Q\otimes1+1\otimes Q$, with the Koszul sign in the
second summand. This construction has a different degree from the
degree-zero exchange operators considered above.
The gauge-theory example makes the relation between the two degrees
explicit \cite[\S\S6.5, 7.2]{DimofteNiuPy2025}.

Let $\mathfrak g$ be a finite-dimensional complex Lie algebra, with
basis $t_a$ and structure constants $[t_a,t_b]=f_{ab}^{\ c}t_c$.
Fix an invariant symmetric form $\kappa$ and a scalar $k\in\mathbb C$.
In this section $k$ denotes the level. The ground field is $\mathbb C$.
Repeated Lie algebra indices are summed.

Write
\[
 h_{\mathfrak g}(x,y)=\tfrac12\operatorname{tr}
                         (\operatorname{ad}x\,\operatorname{ad}y).
\]
For a simple algebra with long roots of square length two,
$h_{\mathfrak g}=h^\vee\kappa$.

\begin{definition}\label{cylev:def:line-algebra}
The graded algebra $\mathcal Y_k$ has even generators $B^a_n$ of
cohomological degree zero and odd generators $c_{a,n}$ of degree one,
for $n\ge0$, with relations
\begin{equation}\label{cylev:eq:line-relations}
 [B^a_n,B^b_m]=f_{ab}^{\ c}B^c_{n+m},\qquad
 [B^a_n,c_{b,m}]=-f_{ac}^{\ b}c_{c,n+m},\qquad
 c_{a,n}c_{b,m}=-c_{b,m}c_{a,n}.
\end{equation}
Define a degree-one derivation on the generators by
\begin{align}
 Qc_{a,n}
   &=\frac12f_{bc}^{\ a}\sum_{m+p=n-1}c_{b,m}c_{c,p},
       \label{cylev:eq:line-ghost}\\
 QB^a_n
   &=nk\kappa_{ab}c_{b,n-1}
     +\frac12f_{ab}^{\ c}\sum_{m+p=n-1}
                   (c_{b,m}B^c_p+B^c_pc_{b,m}).
       \label{cylev:eq:line-even}
\end{align}
Every index in a sum is nonnegative, so the sum is empty for $n=0$.
The weights $\operatorname{wt}(B_n)=n$ and
$\operatorname{wt}(c_n)=n+1$ make $Q$ weight preserving.
The algebra consists of finite polynomial expressions in these modes.
\end{definition}

The symmetric ordering in \eqref{cylev:eq:line-even} determines the
level seen by a cochain model.
Moving every ghost to the left gives
\begin{equation}\label{cylev:eq:line-shift}
 QB^a_n=n(k\kappa+h_{\mathfrak g})_{ab}c_{b,n-1}
       +f_{ab}^{\ c}\sum_{m+p=n-1}c_{b,m}B^c_p.
\end{equation}
Indeed, each of the $n$ commutators contributes
$-\frac12f_{ab}^{\ c}f_{cd}^{\ b}c_{d,n-1}
=(h_{\mathfrak g})_{ad}c_{d,n-1}$.

\begin{proposition}\label{cylev:prop:line-dg}
The derivation $Q$ respects \eqref{cylev:eq:line-relations} and
satisfies $Q^2=0$.
\end{proposition}
\begin{proof}
Put $\mathfrak g_+=\mathfrak g[u]$ and
$\mathfrak g_-=u^{-1}\mathfrak g[u^{-1}]$.
For an invariant form $b$, let $U_b$ be the enveloping algebra of
the central extension
\[
 [xu^p,yu^q]=[x,y]u^{p+q}+p\delta_{p+q,0}b(x,y)\mathbf1,
\]
with the central element specialized to $1$.
Take $b=k\kappa+h_{\mathfrak g}$ and
$M_b=U_b\otimes_{U(\mathfrak g_-)}\mathbb C$.
PBW identifies $M_b$ with $U(\mathfrak g_+)$.

Use restricted Chevalley cochains for $\mathfrak g_-$ on this module.
The ghost dual to $t_au^{-n-1}$ is $c_{a,n}$.
Each cochain is a finite sum of homogeneous energy components.
The energies of $B_n$ and $c_n$ are $n$ and $n+1$, respectively.
Write $\mathfrak g_-^\circ$ for this restricted graded dual.

Choose the ghost sign for which
$dc^x=\frac12c^x([y,z])c^yc^z$ and the action term is
$-\sum_x c^x\rho(x)$.
Here $x,y,z$ range over a homogeneous basis of $\mathfrak g_-$,
$c^x$ is the dual ghost, and $\rho$ is its action on $M_b$.
On $B=t_au^n$, the identity
\[
 [t_bu^{-m-1},B]
 =f_{ba}^{\ c}t_cu^{n-m-1}
       -(m+1)\delta_{m,n-1}b_{ba}
\]
gives exactly \eqref{cylev:eq:line-shift} after acting on the inducing
vector. The ghost differential gives \eqref{cylev:eq:line-ghost}.

This generator calculation extends to products as follows.
Transport left multiplication by $B$ to
$a_B\otimes1+1\otimes L_B$ on
$\Lambda(\mathfrak g_-^\circ)\otimes U(\mathfrak g_+)$, where $a_B$
is the coadjoint derivation for the projected bracket $[B,-]_-$.
The projected Jacobi identity gives
\[
 [d_{\mathrm{CE}},a_B]
       =-\sum_x\varepsilon(c^x)a_{[x,B]_+}.
\]
Here $\varepsilon(c^x)$ is exterior multiplication by $c^x$.
Together with the central term in the displayed bracket, this says
$[d,L_B]=L_{QB}$. For ghost multiplication the same equation is
the exterior Leibniz rule.

Induction on word length now identifies the transported differential
with $Q$ on every PBW monomial.
The multiplication operators already satisfy the relations, so $Q$
respects them. The Chevalley identity $d^2=0$ gives $Q^2=0$.
\end{proof}

Let $T B_n=nB_{n-1}$ and $T c_n=nc_{n-1}$.
This degree-zero derivation commutes with $Q$ and is locally
nilpotent on the polynomial algebra.
Put $\tau_z=e^{zT}$ and let $\epsilon$ kill every generator.
For each generator $a$, define
\begin{equation}\label{cylev:eq:line-spectral}
 \Delta_z(a)=\tau_z(a)\otimes1+1\otimes a,\qquad
 r(z)=\sum_{n,m\ge0}(-1)^n\binom{n+m}{n}
       \frac{c_{a,n}\otimes B^a_m-B^a_n\otimes c_{a,m}}{z^{n+m+1}}.
\end{equation}
The coproduct extends multiplicatively.

The tensor $r$ lies in
$(\mathcal Y_k\otimes\mathcal Y_k)((z^{-1}))$.
This Laurent completion consists of series bounded above in powers
of $z$, with algebraic coefficients. It is larger than algebraic
scalar extension by $\mathbb C((z^{-1}))$.

\begin{theorem}\label{cylev:thm:line-spectral}
The spectral data satisfy
\begin{align}
 Q_\otimes r+r^2&=0,&
 \Delta_z Q&=(Q_\otimes+[r,-])\Delta_z,
          \label{cylev:eq:line-mc}\\
 (\epsilon\otimes1)\Delta_z&=1,&
 (1\otimes\epsilon)\Delta_z&=\tau_z,
          \label{cylev:eq:line-counit}\\
 (T\otimes1+1\otimes T)r&=0,&r(-z)&=P r(z).
          \label{cylev:eq:line-symmetry}
\end{align}
In the first counit identity, $1$ denotes the identity map.
The flip $P$ uses the cohomological sign.
Coassociativity takes the form
\[
 (\Delta_w\otimes1)\Delta_z=(1\otimes\Delta_z)\Delta_{z+w}.
\]
The two iterated twists both equal
$r_{12}(w)+r_{13}(z+w)+r_{23}(z)$.

The shifted kernels are first evaluated on
$0<|w|<|z|<|z+w|$, with sufficiently small auxiliary jet variables.
Their fixed-mode coefficients then continue rationally to
$zw(z+w)\ne0$.
The assertion concerns these kernels and their finite products.
It does not extend translation to every formal Laurent series.
\end{theorem}
\begin{proof}
Translation, the two counits and coassociativity follow on generators
from $\tau_z\tau_w=\tau_{z+w}$.
The binomial coefficients in \eqref{cylev:eq:line-spectral} give
\[
 (\Delta_w\otimes1)r(z)=r_{13}(z+w)+r_{23}(z),\qquad
 (1\otimes\Delta_z)r(z+w)=r_{12}(w)+r_{13}(z+w).
\]
They also give \eqref{cylev:eq:line-symmetry}.
For fixed output modes, the first shifted sum converges normally
when $|w|<|z|$, and the second when $|z|<|z+w|$.
Indeed, their tails are polynomial in the summation index times
the corresponding geometric ratio.
The stated common domain permits both substitutions and their
finite products before rational continuation.

For the two differential identities, write
$B^a(s)=\sum_{n\ge0}B^a_ns^{-n-1}$ and similarly for $c_a(s)$.
Then $r(z)$ is the coefficient of $s^{-1}$ in
$c_a(s)\otimes B^a(s+z)-B^a(s)\otimes c_a(s+z)$, expanded in $s/z$.
The terms of $Q_\otimes r$ containing $k\kappa+h_{\mathfrak g}$
are the residue of a total $s$-derivative and vanish.

Its remaining terms are the residue of
\[
 f_{ab}^{\ c}\left(
 \tfrac12c_ac_b\otimes B^c-c_a\otimes c_bB^c
          -c_bB^c\otimes c_a-\tfrac12B^c\otimes c_ac_b\right).
\]
The first and second factors are evaluated at $s$ and $s+z$.
The commutators \eqref{cylev:eq:line-relations} in $r^2$ give the
negatives of these four terms. This is the residue calculation
of \cite[(7.23)--(7.24)]{DimofteNiuPy2025}.

Both sides of the second identity in \eqref{cylev:eq:line-mc} are
derivations along $\Delta_z$ after subtracting $Q_\otimes\Delta_z$.
On $c_{a,n}$ their common value is
$f_{bc}^{\ a}\sum_{m+p=n-1}\tau_z(c_{b,m})\otimes c_{c,p}$.
On $B^a_n$ it is
\[
 f_{ab}^{\ c}\sum_{m+p=n-1}
       \bigl(\tau_z(c_{b,m})\otimes B^c_p+
                    \tau_z(B^c_p)\otimes c_{b,m}\bigr).
\]
For $\Delta_zQ-Q_\otimes\Delta_z$, these are the cross terms of
the quadratic differential.
For $[r,\Delta_z(-)]$, inserting the commutator relations and summing
the binomial expansion gives the same finite sums
\cite[\S7.2.3]{DimofteNiuPy2025}.
The identity therefore holds on every polynomial.
\end{proof}

For dg left modules $M,N$, the tensor space
$(M\otimes N)((z^{-1}))$ consequently has differential
$D_z=d_\otimes+\rho(r(z))$ and $\mathcal Y_k$-action through $\Delta_z$.
Here $\rho$ denotes the tensor product representation.
Its iterated products use the common triple twist just displayed.
This is the spectral tensor product of the dg line modules.

\subsection{A degree-zero kernel from a translation homotopy}

Assume $k\ne0$ and that $\kappa$ is nondegenerate.
Let $\kappa^{ab}$ be the inverse matrix.
Define a derivation $L$ of degree $-1$ by
\[
 Lc_{a,n}=k^{-1}\kappa^{ab}B^b_n,\qquad LB^a_n=0.
\]
Invariance of $\kappa$ makes these assignments respect all relations.
Put $\Omega=\kappa^{ab}t_a\otimes t_b$.

\begin{proposition}\label{cylev:prop:line-kernel}
The derivation $L$ satisfies $[Q,L]=T$.
The degree-zero tensor $X=(L\otimes1)r$ is
\begin{equation}\label{cylev:eq:line-kernel}
 X(z)=\frac1k\sum_{n,m\ge0}(-1)^n\binom{n+m}{n}
              \frac{\kappa^{ab}B^a_n\otimes B^b_m}{z^{n+m+1}}
      =\frac{\Omega}{k(z+u-v)}.
\end{equation}
In the last expression, $B^a_n$ is read as $t_au^n$ in the first
factor and $t_av^n$ in the second.
On the tensor product of dg modules,
$\nabla_z=\partial_z+\rho(X)$ commutes with $D_z$.
\end{proposition}
\begin{proof}
On $B_n$, the level term of $LQB_n$ is $nB_{n-1}$.
The remaining terms vanish because invariance makes their Lie
coefficient antisymmetric, while the symmetric ordering and the sum
over $(m,p)$ make their even product symmetric.
On $c_n$, the level term of $QLc_n$ is $nc_{n-1}$, and the two
quadratic contributions in $QLc_n+LQc_n$ cancel by the same
invariance. These checks prove $[Q,L]=T$ on the full algebra.

Only the first summand of $r$ survives $L\otimes1$.
Expanding $(z+u-v)^{-1}$ proves \eqref{cylev:eq:line-kernel}.
Apply $L\otimes1$ to the Maurer--Cartan equation.
Since $(L\otimes1)(r^2)=Xr-rX$ and
$(T\otimes1)r=\partial_zr$, the result is
\[
 \partial_zr=Q_\otimes X+[r,X].
\]
This is exactly $[\partial_z+\rho(X),D_z]=0$.
\end{proof}

The inverse in $L$ is the inverse of $k\kappa$.
For a normalized simple algebra, replacing it by
$(k+h^\vee)^{-1}\kappa^{-1}$ gives
$[Q,L]B_1=k(k+h^\vee)^{-1}B_0$, which differs from $TB_1=B_0$
when $h^\vee\ne0$.
Thus the normal-ordering shift in
\eqref{cylev:eq:line-shift} is not a replacement for the level
in the translation homotopy.

A Wilson module is obtained from a $\mathfrak g$-representation
by letting the ghosts and all $B_n$, $n>0$, act by zero.
On two such modules $r$ acts by zero, while the connection becomes
$\partial_z+\rho(\Omega)/(kz)$.
For two defining $\mathfrak{sl}_2$ modules with
$\kappa=\operatorname{tr}_{\mathbb C^2}$,
\[
 \Omega=P-\tfrac12\mathrm{id}.
\]
Its connection residues are $1/(2k)$ and $-3/(2k)$ on the symmetric
and antisymmetric summands.
With affine level $k'=-k-h^\vee$, these are the residues of the
Knizhnik--Zamolodchikov connection in the convention
$d-\Omega\,d\log z/(k'+h^\vee)$.

The ideal generated by the ghosts is stable under $Q$, and its
quotient is $U(\mathfrak g[u])$ with zero differential.
The tensor $r$ maps to zero in this quotient, whereas $X$ survives.
The operation $r\mapsto(L\otimes1)r$ therefore identifies a rational
kernel. It does not identify the dg line algebra with a degree-zero
quantum group.
It also requires $k\ne0$, although the dg spectral structure itself
exists at $k=0$.

\section{The comparisons still to be constructed}
\label{cylev:sec:comparisons}

An algebra isomorphism on states does not determine a field tensor.
Proposition~\ref{cylev:prop:ordinary-obstruction} proves this for the
rational rank-two family, even at level zero.
The dg example adds a second distinction: a differential twist and
its degree-zero connection kernel have different types.
A comparison must specify how these types are related.

Current Ward identities provide a useful intermediate target for
maps from line modules to affine modules.
For smooth affine modules $X_1,X_2,X_3$, such an identity has the form
\[
 xt^m I_z(v_1,v_2)
 =\sum_{j\ge0}\binom mj z^{m-j}I_z(xt^jv_1,v_2)
                         +I_z(v_1,xt^mv_2).
\]
Smoothness means that sufficiently positive modes kill each vector.

Suppose their energy gradings are bounded below with
finite-dimensional energy spaces.
Take $I_z$ with values in
$\prod_e(X_{3,e}\otimes\mathbb C((z^{-1})))$ and require
\[
 \operatorname{pr}_e I_z(X_{1,e_1}\otimes X_{2,e_2})
                   \subset X_{3,e}z^{e-e_1-e_2}.
\]
Then its restricted transpose converts the Ward equation into
linearity for the mode coproduct
\[
 \delta_z(xt^m)=\sum_{j\ge0}\binom{-m}{j}z^{-m-j}
                         xt^{-j}\otimes1+1\otimes xt^m.
\]
The transpose uses the anti-involution $xt^m\mapsto-xt^{-m}$ and
Laurent completion after the algebraic tensor product.
This is a comparison of current actions.
The full intertwining Jacobi identity imposes further equations.

\begin{example}\label{cylev:ex:ward-jacobi}
Let $H_\lambda$ be the Heisenberg vacuum algebra with
$[J_m,J_n]=\lambda m\delta_{m+n,0}$, where $\lambda\ne0$.
For $\alpha\in\mathbb C$, set
$A_\alpha(z)=J(z)+\alpha z^{-1}\mathrm{id}$ and define
\[
 I_\alpha(J_{-n_1}\cdots J_{-n_r}\mathbf1,z)
   ={:}\prod_{i=1}^r
                 \frac{\partial_z^{n_i-1}A_\alpha(z)}{(n_i-1)!}{:}.
\]
The empty product is the identity.
Normal ordering places the negative current modes first.
The scalar summand commutes with every current, so
\[
 [J_m,\partial_z^{n-1}A_\alpha(z)/(n-1)!]
                  =\lambda n\binom mn z^{m-n}\mathrm{id}.
\]

The product rule gives the current Ward identity on every PBW state.
The relation $[L_{-1},J_{-n}]=nJ_{-n-1}$ also gives
$I_\alpha(L_{-1}v,z)=\partial_zI_\alpha(v,z)$.

The full intertwining Jacobi identity, applied with first state
$\mathbf1$, would instead require
$I_\alpha(J_{-1}\mathbf1,z)={:}J(z)I_\alpha(\mathbf1,z){:}=J(z)$.
It fails when $\alpha\ne0$.
Thus the current and derivative equations alone do not construct an
affine intertwining operator.
\end{example}

\begin{problem}\label{cylev:prob:line-field}
Start from a coassociative comparison through an affine bimodule
resolution of the line-module functor.
Choose a derived affine fusion category and coefficient models on
which the comparison computes derived images and admits restricted
transposition.
Does its transpose satisfy the full intertwining Jacobi identity
and the required grading and completion conditions?
If so, determine when the resulting comparison from affine fusion
to the image of the spectral line tensor is invertible.
Retain the unit and normalization on the resolution.
Agreement of Wilson connection residues determines neither this
fusion morphism nor its inverse.
\end{problem}

\subsection{A rational Hall cooperation}

A Hall cooperation can be multiplicative after both its coefficients
and its tensor multiplication are specified.
Set $R_H=\mathbb Q[h_1,h_2,h_3]/(h_1+h_2+h_3)$ and
\[
 N(v)=\prod_{i=1}^3(v+h_i),\qquad
 K_0(v)=N(v)/v,\qquad K_1(v)=vN(v).
\]
For $j=0,1$, give
$\mathcal P_d=R_H[x_1,\ldots,x_d]^{\mathfrak S_d}$
the shuffle product
\[
 (F*_jG)(Z)=
 \sum_{Z=X\sqcup Y}F(X)G(Y)\prod_{x\in X,y\in Y}K_j(y-x),
\]
where the ordered partitions have the prescribed input sizes.

The first kernel has only simple root-diagonal poles, whose paired
residues cancel in the sum. The second is polynomial.
Thus both products are polynomial.
Summing over ordered three-block partitions proves associativity.

For ordered lists $X,Y$ of sizes $d,e$, put
\[
 Q^{(j)}_{d,e}(X,Y)=\prod_{x\in X,y\in Y}K_j(y-x),\qquad
 \Delta^{\mathrm{rat}}_{d,e}F
                      =F(X,Y)/Q^{(j)}_{d,e}(X,Y).
\]
The recipient is the algebra of block-symmetric polynomials localized
at all pair differences and their translates by $\pm h_i$.
No parameter $h_i$ is inverted.
Write $r_H(v)=K_j(-v)/K_j(v)$. The ratio is the same for both $j$.

For bidegrees $(d,e)$ and $(d',e')$, define a product by summing over
$X=X_1\sqcup X_2$, $Y=Y_1\sqcup Y_2$, with these sizes and summand
\begin{equation}\label{cylev:eq:rational-middle}
\begin{split}
 &F(X_1,Y_1)G(X_2,Y_2)
   \prod_{x\in X_1,x'\in X_2}K_j(x'-x)
   \prod_{y\in Y_1,y'\in Y_2}K_j(y'-y)\\
 &\hspace{35mm}\cdot
   \prod_{y\in Y_1,x'\in X_2}r_H(y-x').
\end{split}
\end{equation}
The final factors specify the interchange of the middle two blocks.
They are absent from ordinary tensor multiplication.

\begin{proposition}\label{cylev:prop:rational-hall}
The product \eqref{cylev:eq:rational-middle} is associative and
unital. The sum of the maps $\Delta^{\mathrm{rat}}_{d,e}$ is
multiplicative into this recipient, rationally coassociative and
counital for both charge-zero projections.
\end{proposition}
\begin{proof}
Fix four ordered groups $X_1,Y_1,X_2,Y_2$ in a source product.
Dividing by the kernel between $X_1\sqcup X_2$ and $Y_1\sqcup Y_2$
cancels the factors from $X_1$ to $Y_2$.
The other mixed factors leave
\[
 K_j(x'-y)/K_j(y-x')=r_H(y-x'),
       \qquad y\in Y_1,\ x'\in X_2.
\]
The remaining denominators are those of $\Delta^{\mathrm{rat}}F$
and $\Delta^{\mathrm{rat}}G$.
The factors within the two output lists are exactly those in
\eqref{cylev:eq:rational-middle}.
Summing proves multiplicativity.

For three input pairs, either bracketing has the same factors within
each output list.
It also has the same factor from every earlier right block to every
later left block. Hence the target product is associative.
On three output lists, either iterated cooperation divides by
$Q_{12}^{(j)}Q_{13}^{(j)}Q_{23}^{(j)}$.
The empty lists give the unit and both counits.
All charge sums on a fixed input are finite.
\end{proof}

The construction leaves the polynomial carrier:
\[
 \Delta^{\mathrm{rat}}_{1,1}1=\frac{v}{N(v)}
       \quad(j=0),\qquad
 \Delta^{\mathrm{rat}}_{1,1}1=\frac1{vN(v)}
       \quad(j=1).
\]
It also removes the closed loci at the inverted denominators.
It therefore supplies neither a supported comparison across those
loci nor an ordinary-tensor coproduct on the original polynomial
Hall algebra. The obstruction in
Proposition~\ref{cyjets:prop:no-ordinary-coproduct} retains its scope.

Hall charge now introduces another comparison index.
For a charge monoid $\Gamma$, a Hall cooperation has components
$\Delta_{\gamma_1,\gamma_2}$ from charge
$\gamma_1+\gamma_2$ to a completed tensor of charges
$\gamma_1,\gamma_2$.
The field coproduct $D_{\eta_1,\eta_2}$ instead changes the level
from $\eta_1\eta_2$ to $(\eta_1,\eta_2)$ and takes values in the
covariant field tensor of this chapter.
Even when both structures have a Yangian presentation, these
indices and their coefficient completions remain to be compared.

\begin{problem}\label{cylev:prob:hall-level}
For a specified oriented Hall system and a specified spectral field
family, construct the coefficient and state comparison maps.
Choose a map from Hall charges to levels and prove its compatibility
with their ordered composition.
Construct the tensor comparison that identifies each Hall
cooperation with $D_{\eta_1,\eta_2}$, including the cross contractions,
vacuum maps and expansions at every collision.
Determine the coefficient localization, integral form and completion
on which it is defined and invertible.
\end{problem}

There is also a question internal to the level family.
The formulas above put all $\mathcal V_c$ on a common state module,
but their classical current products have level-dependent central
terms. The translation twist compares two coproduct presentations.
It does not remove that dependence.

One may instead seek an intrinsic category or a category of
coefficient objects in which the different levels admit comparison
functors preserving the relevant tensor.
The first required construction is such a functor together with its
map on the central operation. The ordinary-tensor obstruction fixes
one classical state identification and leaves this broader question
open.
